\chapter{微分博弈理论与 HJI 可达性}
\label{ch:plan-diffgame}

% ════════════════════════════════════════════════════════════════
%  规划部·博弈线地基章：吸收 Tutorial「微分博弈与HJI」。
%  零和微分博弈 -> Isaacs/HJI -> 粘性解 -> 可达性(BRS/BRT/reach-avoid)
%  -> 水平集数值(WENO5/TVD-RK3) -> 维度诅咒三绕行 -> LQ 博弈耦合 Riccati。
%  教学层遵循 v5.0 理论教学规范。
%  跨部去重：CBF/CLF、HJ 可达性、鲁棒/随机 MPC 的最深推导留待【控制部】，
%  本章侧重规划/决策用途（安全滤波、reach-avoid 用于规划、博弈），就地 \pz 标注。
% ════════════════════════════════════════════════════════════════

\begin{practice}[前置自测]
开始之前，先试答下列问题。能流畅作答者，可只速读本章的速查表与符号表；卡住之处，正是本章要为你解决的。这 5 题对应本章反复用到的前置工具。
\begin{enumerate}
  \item 单人有限时域最优控制 $\min_u[\,g(x(T))+\int_t^T L\,\mathrm ds\,]$（受 $\dot x=f(x,u)$ 约束）的值函数 $V(x,t)$ 满足什么偏微分方程？终端条件是什么？
  \item 线性二次调节器（LQR）的微分 Riccati 方程长什么样？给定 Riccati 解 $\mathbf P$，最优反馈增益 $\mathbf K$ 怎么算？
  \item 对一个二元函数 $H(u,v)$，$\min_u\max_v H$ 与 $\max_v\min_u H$ 一般相等吗？什么条件下相等？
  \item 给定标量函数 $V:\mathbb{R}^n\to\mathbb{R}$，集合 $\{x:V(x)\le 0\}$ 叫什么？当 $V$ 取“到某集合的带符号距离”时，这个集合是什么？
  \item Pontryagin 最大值原理（PMP）给出最优轨迹的什么必要条件？协态（costate）在其中扮演什么角色？
\end{enumerate}
\noindent\emph{答案线索}：1$\to$\cref{sec:dg-bridge} 与 \cref{sec:ctrl-lqr}；2$\to$\cref{sec:ctrl-lqr}；3$\to$\cref{sec:dg-isaacs}（Isaacs 条件）；4$\to$\cref{sec:dg-reach}（水平集编码）；5$\to$\cref{sec:dg-isaacs} 的 HJI/PMP 对照。
\end{practice}

\section*{本章目标}
学完本章，你应当能够：
\begin{enumerate}
  \item \textbf{建模}：把一个双智能体对抗问题写成零和微分博弈，写出它的 Isaacs 方程，并判断 Isaacs 条件是否成立（即这个博弈是否“有值”）；
  \item \textbf{推导}：从单人 HJB 出发推广出 HJI 方程，解释为什么 HJI 的解一般不是经典可微解、必须是黏性解，并用消失黏性法定出黏性解的次/超解条件；
  \item \textbf{关联可达性}：把“安全验证”转化为可达性计算——分清后向可达集（BRS）、后向可达管（BRT）、reach--avoid 集三者，写出各自的 HJI 方程，并由玩家角色（控制 vs 干扰、存在 vs 任意）唯一确定哈密顿量里 $\min/\max$ 的取法；
  \item \textbf{理解数值方法}：说清为什么要用水平集方法把“是否到达”这种布尔问题转成连续值函数的演化，以及网格离散带来的维度诅咒从何而来，掌握 WENO5 + TVD-RK3 + CFL 的设计动机；
  \item \textbf{推导耦合 Riccati}：写出 $N$ 人 LQ 博弈的反馈 Nash 条件与耦合 Riccati 方程组，理解它为什么“耦合”、为什么不保证有解——这是后续实时博弈求解器的直接前置；
  \item \textbf{掌握绕维度诅咒的三条路}：状态解耦/STP、FaSTrack、神经近似（DeepReach），并能按“精确性 vs 可扩展性”权衡选型；
  \item \textbf{落地经典问题}：解释 Homicidal Chauffeur 与 Air3D 这两个追逃博弈基准，以及为什么要在相对坐标下建模。
\end{enumerate}

\subsection*{本章知识导航}
本章是一条“\emph{从对抗交互下的“最优”定义，一路下降到可算近似}”的主线。树根是一个问题——对抗交互下“最优”是什么；树干是从理论方程到可算近似的下降链；树枝是数值方法与经典问题。结构如下：

\begin{center}
\begin{forest} navtreeLR
[{博弈 = 最优控制\\ +“对手也在最优”}
  [{Isaacs / HJI\\ $\min_u\max_v$、鞍点}
    [{黏性解\\ 消失黏性、唯一性}]
    [{可达性\\ BRS / BRT / reach--avoid}
      [{水平集数值\\ WENO5 + TVD-RK3}
        [{维度诅咒\\ 解耦·FaSTrack·DeepReach}]
        [{LQ 博弈\\ 耦合 Riccati}
          [{追逃经典\\ Chauffeur · Air3D}]
        ]
      ]
    ]
  ]
]
\end{forest}
\end{center}

\noindent\textbf{推荐路径}：\cref{sec:dg-why}（为什么需要博弈视角）定下世界观，无公式但奠定整章动机；\cref{sec:dg-isaacs,sec:dg-viscosity,sec:dg-reach}（Isaacs/HJI、黏性解、可达性）是本章\emph{最核心、最难}的三角，建立“博弈—方程—弱解—可达集”的逻辑链；\cref{sec:dg-numerics}（水平集数值）面向工程实现，初读可只看动机与 WENO/CFL 结论；\cref{sec:dg-lq}（LQ 博弈与耦合 Riccati）是\emph{本章最关键}的一节，是后续实时博弈求解器的直接前置，整章若只精读一节就是它；\cref{sec:dg-curse}（维度诅咒三绕行）面向选型；\cref{sec:dg-chase}（追逃经典）把全部理论落到两个标准基准。难度集中在 \cref{sec:dg-isaacs,sec:dg-viscosity,sec:dg-reach,sec:dg-lq}；时间紧张时 \cref{sec:dg-numerics}（数值细节）与 \cref{sec:dg-chase}（经典问题）可先速读。

\subsection*{前置知识桥接}
\label{sec:dg-bridge}
本章建立在单人最优控制之上，这里把会用到的三块前置激活一下，读者不必翻回前置章也能跟上。

\paragraph{HJB 方程（单人，最小化）。}
考虑 $\min_u[\,g(x(T))+\int_t^T L(x,u)\,\mathrm ds\,]$，受 $\dot x=f(x,u)$ 约束。其值函数 $V(x,t)$（从 $(x,t)$ 出发的最优剩余代价）满足 Hamilton--Jacobi--Bellman（HJB）方程
\begin{equation}
  -\frac{\partial V}{\partial t}=\min_{u}\Big[\nabla_x V\cdot f(x,u)+L(x,u)\Big],\qquad V(x,T)=g(x).
  \label{eq:dg-hjb}
\end{equation}
直觉是：最优剩余代价随时间往前推进的变化率，等于“当前这一步最优地花掉的瞬时代价 + 状态移动带来的未来代价变化”。这里 $\nabla_x V$ 就是 PMP 里的协态 $\lambda$——值函数对状态的梯度即“影子价格”。本章的 HJI 方程，就是把这个 $\min_u$ 换成博弈的 $\min_u\max_v$。\cref{ch:control} 已系统建立单人最优控制与 LQR 框架，本章只取其结论用作桥接，不重复推导。

\paragraph{LQR 与 Riccati。}
线性系统 $\dot{\mathbf x}=\mathbf A\mathbf x+\mathbf B\mathbf u$、二次代价 $\int(\mathbf x^\top\mathbf Q\mathbf x+\mathbf u^\top\mathbf R\mathbf u)\,\mathrm ds$ 的最优反馈是 $\mathbf u^*=-\mathbf K\mathbf x$，增益 $\mathbf K=\mathbf R^{-1}\mathbf B^\top\mathbf P$，其中 $\mathbf P$ 解 Riccati 方程
\begin{equation}
  \dot{\mathbf P}=-\mathbf A^\top\mathbf P-\mathbf P\mathbf A+\mathbf P\mathbf B\mathbf R^{-1}\mathbf B^\top\mathbf P-\mathbf Q\quad(\text{终端给定}).
  \label{eq:dg-riccati}
\end{equation}
本章 \cref{sec:dg-lq} 会看到：$N$ 人 LQ 博弈把这一条 Riccati 变成 $N$ 条相互耦合的 Riccati——这就是从“单人最优”到“多人均衡”在线性世界里的全部数学增量。Riccati 解矩阵 $\mathbf P$ 与本书估计章的协方差同字母但语义相反（一个是代价权重，一个是不确定度），沿用 \cref{ch:control} 的约定，本章 $\mathbf P$ 一律指代价权重。

\paragraph{带符号距离与水平集。}
给定闭集 $\mathcal{T}$，其\emph{带符号距离函数} $g(x)$ 满足：$x\in\mathcal{T}\Leftrightarrow g(x)\le 0$，且 $|g(x)|$ 是 $x$ 到边界 $\partial\mathcal{T}$ 的距离。于是“集合”被编码成“函数的零下水平集” $\{x:g(x)\le 0\}$，集合的演化变成函数的演化。这是 \cref{sec:dg-reach,sec:dg-numerics} 把布尔问题（到没到达）转成连续问题（值函数）的关键技巧。

\subsection*{如果跳过本章会怎样}
\begin{itemize}
  \item 在做交互式导航（路口、人流、并线）时，你会沿用 \cref{ch:planning} 的“预测—然后—规划”管线，把对向车/行人当固定移动障碍，却\emph{不理解}为什么会陷入 frozen robot（僵住的机器人）——因为你切断了“我动一下、对方会让一下”的反馈环；
  \item 在做安全验证时，你会想用一个集合刻画“危险状态”，却\emph{分不清}后向可达集（BRS）与后向可达管（BRT），把只看终点的方程用到本该看整个时间窗的安全分析上，得到偏乐观、不安全的结论；
  \item 在读实时博弈求解器（iLQGames 一类）的代码时，你会看到一组“耦合 Riccati 方程”却\emph{不知道}它就是本章的 LQ 博弈反馈 Nash 条件，只能照抄而不会调试。
\end{itemize}

\begin{table}[H]\centering\small
\caption{本章阅读方式建议}
\label{tab:dg-reading}
\begin{tabular}{lll}
\toprule
阅读方式 & 时间 & 适合谁 \\
\midrule
精读（含手推 Isaacs/Riccati + 跑一个 Air3D） & 12--16 小时 & 首次系统学习、要做博弈/可达性研究 \\
速读（理解概念脉络，跳过数值细节与手推） & 4--6 小时 & 有最优控制基础、想抓主线 \\
速查（只看小结的符号表 + 速查表） & 20 分钟 & 写代码、查公式时回来 \\
\bottomrule
\end{tabular}
\end{table}

\begin{note}
\textbf{本章记号与约定.}\;
本章处于规划/决策语境，沿用该领域自然记号，并与 \cref{ch:planning,ch:control} 一致。状态 $x$（或 $\mathbf x$）、动力学 $f(x,u,v)$；两玩家的控制在博弈一般理论里记 $u,v$，在可达性一节沿用文献记号 $a$（控制）、$b$（干扰）。值函数 $V(x,t)$，其对状态的梯度即\emph{协态} $\lambda=\nabla_x V$，对时间的偏导在可达性方程中沿用文献记号 $D_t V=\partial V/\partial t$。哈密顿量 $H(x,\lambda)$。目标集 $\mathcal{T}$、其带符号距离 $g$；危险集 $\mathcal{A}$、其带符号距离 $c$（均用“集合内 $\le0$”的标准约定）。LQ 博弈沿用控制记号 $\mathbf A,\mathbf B_i,\mathbf Q_i,\mathbf R_{ii},\mathbf P_i,\mathbf K_i$（\cref{ch:control} 约定，$\mathbf P_i$ 为代价权重非协方差）。\textbf{字母隔离声明}：本章 $\lambda$ 指协态（非特征值/波长）、$a,b$ 指控制/干扰（非加速度）、$c$ 指危险集距离（非光速/代价），均仅本章局部生效。Air3D 用相对位姿在李群上叙述时与 \cref{ch:lie} 的 $\SO(2)/\SE(2)$ 约定一致；本章不直接用 $\SO(3)$ 记号，但凡涉及刚体相对位姿的追逃，群/代数记号一律用 preamble 宏 $\SO$、$\SE$、$\so$、$\se$。
\end{note}

% ════════════════════════════════════════════════════════════════
\section{从最优控制到微分博弈：为什么需要博弈视角}
\label{sec:dg-why}

\paragraph{动机：一个十字路口的两难。}
设想一个无保护左转。你的车要左转穿过对向直行车道，对向有一辆车正在接近。规划器需要决定：抢在对方前面转过去，还是让对方先过。把这个场景交给一个标准的“预测—然后—规划”管线（\cref{ch:planning} 的级联范式），它会先用预测模块预测对向车未来几秒的轨迹，把这条轨迹当作一个\emph{移动的障碍物}；然后在“避开这个障碍物”的约束下，为自己求一条最优轨迹。问题出在第一步：预测模块预测对向车轨迹时，\textbf{假设对向车的行为与你无关}——它没建模“如果你果断把车头探进路口，对向车很可能会减速让你”。换句话说，对向车的真实轨迹\emph{取决于你的行为}，而你的规划又取决于对它的预测。这是一个鸡生蛋的循环，“预测—然后—规划”把它粗暴地切断：先固定一方，再优化另一方。

\paragraph{反面：被切断的反馈环会怎样。}
把这个循环切断，会催生两种典型病灶。

\textbf{病灶一：frozen robot（僵住的机器人）。}\;若预测模块给出“对向车会一直匀速直行”（因为它没建模交互），那么在你看来，整条对向车道在未来几秒内都被占满，没有任何安全间隙。于是规划器得出唯一“安全”的结论：原地不动，等对方过完。可现实里只要你稍微往前探，对方就会让——但模型里没有这个机制，于是你永远在等，对方也因为你不动而继续直行，僵局形成。在密集人流中导航的机器人最容易犯这个病：把每个行人当成不可协商的移动障碍，结果寸步难行。量化地说，当行人密度超过某阈值，“把所有人当固定障碍”的规划器会发现可行集为空、完全停摆——而人类在同样人流里照样穿行，靠的正是“我动一下、对方会让一下”的隐式协商。

\textbf{病灶二：预测与规划的不一致。}\;若预测器乐观地认为“对方会让我”，而对方其实是个不让的激进司机，规划器就会做出危险抢行；反过来若预测器悲观，又退化成 frozen robot。根子是同一个：\textbf{预测把对手的策略当成了与己无关的外生输入}，而它本该是内生的、对你的行为做出反应的。

\begin{insight}
交互场景里真正缺的不是“更准的预测器”，而是承认对手是一个\textbf{有目标的决策者}——它的轨迹不是一条要去拟合的曲线，而是它自己优化问题的解。一旦接受这一点，“预测”就不再是独立模块：对手的“预测轨迹”应当是博弈均衡中对手的最优响应。这就是“预测即均衡”的思想。
\end{insight}

\paragraph{核心思想：三次认知跨越。}
把上面的直觉提炼一下，从单智能体规控走到博弈规控，要完成三次认知跨越。这三句话是整条博弈线的精神主轴。

\textbf{跨越一：从“最优控制”到“均衡”。}\;单智能体规控的核心是 $\min_u J(x,u)$——凸时存在唯一全局最优（\cref{ch:control}）。博弈规控没有“最优”的概念，只有“\textbf{均衡}”——每个人在其他人策略固定时不能单方面改善。把单智能体和博弈并排，差异一目了然（\cref{tab:dg-vs}）。这张表里最容易被低估、却最重要的一行是“唯一性”：在十字路口，“你让我先过”和“我让你先过”\emph{都是}合法的纳什均衡（双方都没有单方面改善的动机），数学不会告诉你该选哪个——\textbf{均衡选择是一个设计决策，而非一个数学结果}。这是从单智能体到博弈最大的认知转变：你不再是在求一个确定的最优，而是在一族均衡里，用额外准则（社会习俗、礼让规则、谁先公示意图）去挑一个。本章 \cref{sec:dg-isaacs} 的 Isaacs 方程，正是零和博弈里刻画均衡的工具。

\begin{table}[H]\centering\small
\caption{单智能体最优控制 vs 微分博弈}
\label{tab:dg-vs}
\begin{tabular}{p{0.20\textwidth} p{0.34\textwidth} p{0.34\textwidth}}
\toprule
维度 & 单智能体最优控制 & 微分博弈 \\
\midrule
目标 & $\min_u J(x,u)$ & 每人有各自的 $J_i$ \\
“最好”的含义 & 全局最优（凸时唯一） & \textbf{均衡}：他人固定时无法单方面改善 \\
解的存在性 & 通常存在 & 纯策略下可能不存在 \\
解的唯一性 & 凸时唯一 & 通常\textbf{不唯一}（选哪个是设计决策） \\
求解 & 一个优化问题 & 耦合的多个优化问题（求不动点） \\
\bottomrule
\end{tabular}
\end{table}

\textbf{跨越二：从“预测-然后-规划”到“预测即均衡”。}\;传统管线先用预测器预测他人轨迹、再把预测当固定障碍做规划——这正是上面 frozen robot 的根源。博弈规划把预测和规划统一：\textbf{他人的“预测”就是博弈均衡中他人的最优响应}，不需要独立的预测模块。ego 的“规划”是均衡里 ego 的最优策略，他人的“预测”是同一个均衡里他人的最优策略，两者同时求解、自洽。

\textbf{跨越三：从“我的代价函数已知”到“对手的代价函数需要推断”。}\;单智能体规控的 $J(x,u)$ 由设计者确定；博弈规控中，\textbf{对手的代价函数 $J_i$ 往往未知}——需要从观测轨迹反推（逆博弈/逆强化学习）。这把优化问题从“给定目标求最优”升级为“先推断目标、再求均衡”的双层结构。

\begin{note}
\textbf{不是 X 而是 Y.}\;博弈规划\textbf{不是}“更复杂版本的最优控制”，而是一类\textbf{结构不同}的问题。最优控制求一个点（最优解），博弈求一个不动点（均衡）；最优控制问“怎样最好”，博弈问“在别人也最好的前提下，什么是稳定的”。把博弈当成“多约束的最优控制”来套求解器，会在均衡不唯一、不存在时栽跟头。
\end{note}

\paragraph{历史：Isaacs 与微分博弈的起源。}
把“对手也在最优决策”这件事数学化的，是 Rufus Isaacs。1950 年代他在兰德公司（RAND Corporation）研究空战与追逃问题，提出了“微分博弈（differential game）”框架，并在 1965 年出版专著 \emph{Differential Games}（Wiley）系统化了它\cite{isaacs1965differential}。Isaacs 的标志性问题——\cref{sec:dg-chase} 会细讲的 Homicidal Chauffeur（“开车的杀手”）——就是一个追逃博弈：转弯半径受限的追捕者去抓机动灵活的逃跑者，双方在连续时间里实时最优地决策。Isaacs 的关键贡献不只是提出问题，而是给出了刻画双方最优策略的方程（今天叫 Isaacs 方程/HJI 方程，\cref{sec:dg-isaacs}），并发现了一个深刻的对称性条件（Isaacs 条件，决定博弈是否“有值”）。他还区分了两类博弈，后面会反复用到：

\begin{itemize}
  \item \textbf{程度博弈（game of degree）}：结果是一个连续数值（如最终被抓住时的距离、追上所需时间），双方优化这个数值；
  \item \textbf{类型博弈（game of kind）}：结果是一个布尔判断（追上了没有、撞了没有），双方争的是定性胜负。
\end{itemize}

这套理论后来经 Ba\c{s}ar 与 Olsder 在 \emph{Dynamic Noncooperative Game Theory}（SIAM Classics 1999）中系统整理为动态/微分博弈的统一框架\cite{basar1999dynamic}，又经 Evans 与 Souganidis（1984）从黏性解理论给出严格的数学基础\cite{evans1984differential}。\cref{sec:dg-reach} 还会看到一个漂亮的转化：水平集方法能把“类型博弈”（到没到达）化成“程度博弈”（值函数的零下水平集），从而可数值求解——这是整个可达性领域的奠基技巧。

\paragraph{零和 vs 一般和：本章的边界。}
博弈分两大类，本章重心在零和。\textbf{零和博弈}：一方所得即另一方所失，$J_1=-J_2$；追逃、安全验证（你 vs 最坏情况干扰）属此类，它有最干净的理论（HJI 方程，\cref{sec:dg-isaacs,sec:dg-viscosity,sec:dg-reach}）。\textbf{一般和博弈}：各方目标不完全对立也不完全一致（如交互式驾驶——大家都想快、也都想不撞），理论更复杂、均衡更易不唯一。本章先在零和的干净世界里把 HJI 方程、可达性、耦合 Riccati 这些工具打磨好。一个提醒：\cref{sec:dg-lq} 的 LQ 博弈是个例外——它本身就是一般和（每人有自己的 $\mathbf Q_i,\mathbf R_i$），放在本章是因为它的耦合 Riccati 是后续实时博弈求解器的直接数学前置，与零和理论无缝衔接。下表给出博弈在机器人学中的“出没地图”，并标出每类问题更适合零和还是一般和——这个区分直接决定用哪套工具。

\begin{table}[H]\centering\small
\caption{博弈在机器人学中的出没地图}
\label{tab:dg-map}
\begin{tabular}{p{0.30\textwidth} p{0.22\textwidth} p{0.20\textwidth} p{0.18\textwidth}}
\toprule
场景 & 玩家 & 博弈类型 & 本章对应 \\
\midrule
无人机空战 / 拦截 & 拦截方 vs 入侵方 & 零和（纯对抗） & \cref{sec:dg-chase,sec:dg-isaacs} \\
安全验证（系统 vs 最坏扰动） & 控制 vs 干扰 & 零和（鲁棒视角） & \cref{sec:dg-reach} \\
无保护左转 / 并线汇入 & 车 vs 车 & 一般和 & \cref{sec:dg-lq}（→ iLQGames） \\
人机协作 / 自驾与人类 & 机器 vs 人 & 一般和 + 层级 & \cref{sec:dg-lq} 后续线 \\
多机器人协调 / 覆盖 & 多机器人 & 一般和（势博弈） & 后续博弈线 \\
\bottomrule
\end{tabular}
\end{table}

\begin{note}
\textbf{多视角理解（零和 $\leftrightarrow$ 鲁棒控制）.}\;“安全验证”这一行值得单独点出：它表面上不是博弈，但把“最坏情况干扰”想成一个对抗玩家，它就变成了纯零和博弈。相似之处：干扰像一个处处与你作对的对手，你求“在它最坏发挥下仍安全”的控制——这正是 $\min_u\max_v$。不同之处：真实干扰并没有“意图”，它只是恰好取到最坏值；把它拟人化为玩家只是为了用博弈工具，不要因此以为风、噪声真的在“算计”你。这条把零和博弈和鲁棒控制接起来的视角，是 \cref{sec:dg-reach} 可达性的共同出发点。\pz{跨部去重：鲁棒/随机 MPC 的“最坏情况干扰”深推导留待控制部 MPC 章，此处仅取其博弈视角。}
\end{note}

\begin{pitfall}[概念误区：把博弈当成“加了对手约束的最优控制”]
新手常想“那我把对手的最优轨迹当约束塞进我的优化不就行了？这和避障没区别”。\textbf{后果}：在均衡不唯一时，求解器随机停在某个均衡上，行为不可预测；在 frozen robot 场景下导致僵死或危险抢行。\textbf{根因}：对手的轨迹不是外生约束，而是它自己优化问题的解，且这个解依赖于你的策略——把它当固定约束，等于人为切断了交互的反馈环。\textbf{正解}：把双方的优化问题联立，求\emph{均衡}（不动点）而非单方最优。检验法：问“如果我换个策略，对手的‘最优轨迹’会不会变？”——会变，就说明它不该被当固定约束。
\end{pitfall}

\begin{pitfall}[思维陷阱：默认博弈均衡唯一、必然存在]
把单智能体“凸问题有唯一全局最优”的直觉直接搬过来，假设博弈也有唯一确定的解。\textbf{后果}：系统在多均衡场景下行为漂移（同样输入时而让行时而抢行）；在纯策略均衡不存在的博弈里求解器不收敛却找不到原因。\textbf{根因}：纳什均衡是不动点，不动点可能多个、可能没有（纯策略下）、可能不稳定；存在唯一性需要额外条件（如势博弈、严格凹凸），不是默认。\textbf{正解}：先判断博弈类型与均衡结构，接受“均衡选择是设计决策”，显式引入选择准则（礼让规则、初始化偏好等）。
\end{pitfall}

\begin{practice}
\begin{enumerate}
  \item （开放分析）找一个你亲身遇到过的“预测—然后—规划”会失败的真实交互场景（路口、电梯口、人流、并线），写出：（a）若把对手当固定障碍，规划器会得出什么结论；（b）这个结论为什么不合理；（c）“承认对手在最优决策”会如何改变结论。把它对应到三次认知跨越中的哪一次。
  \item （概念）用自己的话解释为什么“均衡选择是设计决策而非数学结果”，举一个均衡不唯一的日常例子（不限于交通），并说明人类社会用什么“额外准则”挑选均衡。
  \item （辨析）下面四个问题分别更适合零和、一般和，还是根本不需要博弈？说明理由：(i) 无人机空战拦截；(ii) 两车并线汇入；(iii) 机器人在空旷房间走向充电桩；(iv) 扫地机器人与家里的猫“抢”同一条走道。
\end{enumerate}
\end{practice}

% ════════════════════════════════════════════════════════════════
\section{零和微分博弈与 Isaacs 方程}
\label{sec:dg-isaacs}

\paragraph{动机：追逃博弈里“同一个代价，一方想小一方想大”。}
零和博弈最纯粹的形态是追逃。设状态 $x$（比如追捕者与逃跑者的相对位置）按 $\dot x=f(x,u,v)$ 演化，$u$ 是玩家 1（追捕者）的控制、$v$ 是玩家 2（逃跑者）的控制。定义代价
\begin{equation}
  J(x,t,u(\cdot),v(\cdot))=g\big(x(T)\big)+\int_t^T L\big(x(s),u(s),v(s)\big)\,\mathrm ds,
  \label{eq:dg-cost}
\end{equation}
比如 $g(x(T))$ 取终端时刻两者的距离。追捕者想让这个距离\emph{小}（$\min_u$），逃跑者想让它\emph{大}（$\max_v$）。同一个 $J$，一方求最小、一方求最大——这就是零和。问题是：双方同时连续地决策，谁也不能预知对方下一刻的动作，最优策略到底长什么样？

\paragraph{反面：天真地“我先优化、你再优化”会丢掉什么。}
一个诱人的错误是把它拆成两步：先固定 $v$ 让 $u$ 求最优，再固定这个 $u$ 让 $v$ 求最优。但这会丢掉博弈的核心——\textbf{同时性}与\textbf{信息结构}。若让 $u$ 先承诺一条完整轨迹、$v$ 再针对它最优响应，那 $v$ 占了全部便宜（它知道 $u$ 的全部计划）；反过来则 $u$ 占便宜。这两种拆法给出不同结果，且都不是“双方同时实时博弈”的真相。真相介于两者之间，由\textbf{值的上下界是否相等}来界定——这正是 Isaacs 条件要回答的。

\subsection{先看最简单的情形：2×2 矩阵博弈与鞍点}
\label{sec:dg-matrix}
连续时间的 Isaacs 条件听起来抽象，但它的内核在一个\emph{有限的} 2×2 矩阵博弈里就能看清。先把这个台阶踩稳。设两个玩家各有两个动作：玩家 1 选行、玩家 2 选列，收益写成矩阵 $M$，$M_{ij}$ 是“玩家 1 付给玩家 2”的数额——零和，故\textbf{玩家 1 想让它小、玩家 2 想让它大}。和连续情形一样，先后手带来差异，于是有两个值：

\begin{itemize}
  \item \textbf{下值（maxmin）}：玩家 2 先“承诺”选哪列，玩家 1 看到后挑该列最小的数（行最小）；玩家 2 则挑能让这个行最小\emph{最大}的列：$\underline V=\max_j\min_i M_{ij}$。
  \item \textbf{上值（minmax）}：玩家 1 先承诺选哪行，玩家 2 挑该行最大的数（列最大）；玩家 1 挑能让这个列最大\emph{最小}的行：$\overline V=\min_i\max_j M_{ij}$。
\end{itemize}

\begin{example}[有鞍点 vs 无鞍点]
\label{ex:dg-saddle}
\textbf{例一（有鞍点）。}\;取 $M=\big[\begin{smallmatrix}2&3\\1&4\end{smallmatrix}\big]$。每行的最大是 $3,4$，故 $\overline V=\min(3,4)=3$（玩家 1 选第 1 行）；每列的最小是 $1,3$，故 $\underline V=\max(1,3)=3$（玩家 2 选第 2 列）。两者相等：$\overline V=\underline V=3$。格子 $(r_1,c_2)=3$ 是一个\textbf{鞍点}——它同时是所在行的最大、所在列的最小，谁都不想单方面偏离。这就是博弈“有值”的离散版：先后手无所谓，纯策略下存在稳定均衡。

\textbf{例二（无鞍点）。}\;取 $M=\big[\begin{smallmatrix}3&1\\2&4\end{smallmatrix}\big]$。行最大 $3,4$ 故 $\overline V=3$；列最小 $2,1$ 故 $\underline V=\max(2,1)=2$。这次 $\overline V=3>\underline V=2$，\textbf{没有纯策略鞍点}：谁先暴露策略谁就吃亏，差额 $3-2=1$ 就是“信息优势”的价值。要让博弈重新“有值”，得允许\textbf{混合策略}（按概率随机选行/列）——von Neumann 的 minimax 定理保证：任何有限零和博弈在混合策略下 $\overline V=\underline V$，总有鞍点。
\end{example}

\begin{note}
\textbf{对比性思维（矩阵博弈 $\leftrightarrow$ 微分博弈的 Isaacs 条件）.}\;例一（有鞍点、$\overline V=\underline V$）正是 Isaacs 条件\emph{成立}的离散缩影；例二（无鞍点）则是 Isaacs 条件\emph{失效}的缩影。微分博弈里那个 $\min_u\max_v$ 与 $\max_v\min_u$ 是否相等，就是在问“逐点的哈密顿量这个小博弈有没有鞍点”——和矩阵博弈是同一个问题，只是把有限矩阵换成关于连续动作 $(u,v)$ 的函数。不同之处：矩阵博弈靠混合策略总能恢复鞍点，而微分博弈的 Isaacs 条件讨论的是\emph{纯反馈策略}下的鞍点，依赖哈密顿量的可分离/凸凹结构、不总成立。记住这个台阶：以后看到 $\min\max$ 与 $\max\min$，先在脑子里画一个 2×2 矩阵，问“有没有一个格子同时是行最大、列最小”。
\end{note}

\subsection{上值、下值与信息结构}
\label{sec:dg-info}
严格定义博弈的解，要先固定\textbf{信息结构}——谁能看到谁的什么。一种标准做法（也是 \cref{sec:dg-reach} 可达性里用的）是\textbf{非预期策略（non-anticipative strategy）}，即 Elliott--Kalton 策略：让一方可以根据对方\emph{到当前时刻为止}的控制来决定自己的动作，但不能预知未来。\textbf{下值}把这个瞬时信息优势给\emph{最小化方} $u$：设 $u$ 的非预期策略 $\alpha$ 是一个从“$v$ 的历史”到“$u$ 的动作”的映射，满足因果性（若两个 $v$ 在 $[t,s]$ 上一致，则 $\alpha$ 对它们的响应在 $[t,s]$ 上也一致），$v$ 则取开环控制。于是
\begin{equation}
  V^-(x,t)=\inf_{\alpha}\ \sup_{v(\cdot)}\ J\big(x,t,\alpha[v](\cdot),v(\cdot)\big),
  \label{eq:dg-lower}
\end{equation}
其中 $\alpha$ 取遍 $u$ 的非预期策略。对每个固定的开环 $v$，$\alpha$ 能逐点取到 $\min_u$，故 $V^-$ 正对应逐点哈密顿量的 $\max_v\min_u$（先动者 $v$ 吃亏、$u$ 见招拆招）。对称地，把信息优势改给\emph{最大化方} $v$（$v$ 用非预期策略、$u$ 开环）定义\textbf{上值} $V^+=\sup_{\beta}\inf_{u(\cdot)}J(x,t,u(\cdot),\beta[u](\cdot))$，它对应 $\min_u\max_v$。直觉上，上、下值之别就是“让谁占信息便宜”。一般地有 $V^-\le V^+$（\cref{prop:dg-maxmin}）。

\begin{proposition}[先动者吃亏：$\max\min\le\min\max$]
\label{prop:dg-maxmin}
对任意二元函数 $H(u,v)$，恒有
\begin{equation}
  \max_{v}\min_{u}H(u,v)\ \le\ \min_{u}\max_{v}H(u,v).
  \label{eq:dg-maxmin}
\end{equation}
\end{proposition}
\begin{proof}
固定任意 $v_0$：$\min_u H(u,v_0)$ 是 $H(\cdot,v_0)$ 在所有 $u$ 上的最小值，按最小值定义不大于任何具体的 $H(u,v_0)$，即 $\min_u H(u,v_0)\le H(u,v_0)$ 对一切 $u$ 成立。两边对 $v_0$ 取最大（左边仍依赖 $v_0$）：$\max_v\min_u H(u,v)\le\max_v H(u,v)$ 对任意固定 $u$ 成立。右边对一切 $u$ 成立，故它也不小于“挑最好的 $u$”，对左边再取 $\min_u$ 不破坏不等式，即得 \cref{eq:dg-maxmin}。
\end{proof}

\begin{insight}
\cref{eq:dg-maxmin} 的物理含义是\textbf{信息优势永远非负}。内层先动的一方（被 $\min/\max$ 套在里面、要先确定）处于劣势，因为外层的对手能看到它的选择再回应。谁后动、谁能“见招拆招”，谁就不吃亏。等号成立——即先后手无所谓——是一个特殊而美好的情形，由 Isaacs 条件刻画。
\end{insight}

\subsection{动态规划与 HJI 方程的推导}
\label{sec:dg-hji-derive}
下面把刻画博弈的值 $V$（Isaacs 条件下上、下值相等，统一记为 $V$）的偏微分方程一步步推出来。核心工具是动态规划原理：把时间窗 $[t,T]$ 切成很短的第一段 $[t,t+\delta]$ 与剩下的 $[t+\delta,T]$。第一段里双方博弈一小步，第二段的剩余博弈递归地由值函数 $V$ 自己描述，给出动态规划递归式（以下值 $V^-$ 的非预期策略形式写出，$\alpha$ 为 $u$ 的非预期策略）：
\begin{equation}
  V(x,t)=\sup_{v}\inf_{\alpha}\left[\int_t^{t+\delta}L\big(x(s),\alpha[v],v\big)\,\mathrm ds+V\big(x(t+\delta),\,t+\delta\big)\right].
  \label{eq:dg-dpp}
\end{equation}
直觉是：从 $(x,t)$ 出发的对抗最优值，等于“第一小段累积的瞬时代价”加上“走到新状态后剩余博弈的值”，双方在第一段把这个总和博弈到鞍点。这一步的合法性正是 \cref{sec:dg-viscosity} 的黏性解理论所保证的——这里先用它推方程，严格性回头补。

\begin{derivation}[从动态规划到 HJI 方程]
\label{deriv:dg-hji}
\textbf{第一步，处理第二段。}\;当 $\delta$ 很小时状态只移动一点点，$x(t+\delta)\approx x+\delta f(x,u,v)$，对 $V$ 在 $(x,t)$ 处做一阶 Taylor 展开：
\[
  V\big(x(t+\delta),\,t+\delta\big)\approx V(x,t)+\delta\,\frac{\partial V}{\partial t}+\delta\,\nabla_x V\cdot f(x,u,v).
\]
\textbf{第二步，处理第一段的积分。}\;区间长 $\delta$ 很短，被积函数近似为常值，$\int_t^{t+\delta}L\,\mathrm ds\approx\delta\,L(x,u,v)$。

\textbf{第三步，代回 \cref{eq:dg-dpp}。}\;等式两边都出现的 $V(x,t)$ 相互抵消，剩下
\[
  0=\sup_{v}\inf_{\alpha}\Big[\delta\,L(x,u,v)+\delta\,\frac{\partial V}{\partial t}+\delta\,\nabla_x V\cdot f(x,u,v)\Big].
\]
\textbf{第四步，两边除以 $\delta$ 并令 $\delta\to 0$。}\;当步长趋于零，第一段的非预期策略博弈退化为关于瞬时动作的逐点静态博弈，$\sup_v\inf_\alpha$ 就变成逐点的 $\max_v\min_u$；在 Isaacs 条件下它与 $\min_u\max_v$ 相等（内部博弈有鞍点），故可写成习惯的 $\min_u\max_v$；而 $\partial V/\partial t$ 不含 $u,v$，可提到博弈外面：
\[
  0=\frac{\partial V}{\partial t}+\min_u\max_v\Big[\nabla_x V\cdot f(x,u,v)+L(x,u,v)\Big].
\]
移项并补上终端条件（$t=T$ 时已无剩余博弈，值即终端代价），即得 HJI 方程。
\end{derivation}

\begin{theorem}[Hamilton--Jacobi--Isaacs（HJI）方程]
\label{thm:dg-hji}
在 Isaacs 条件下，零和微分博弈的值 $V$（此时上、下值相等）满足
\begin{equation}
  -\frac{\partial V}{\partial t}=\min_{u}\max_{v}\Big[\nabla_x V\cdot f(x,u,v)+L(x,u,v)\Big],\qquad V(x,T)=g(x).
  \label{eq:dg-hji}
\end{equation}
（若 Isaacs 条件不成立，则上值 $V^+$ 满足此式、下值 $V^-$ 满足把 $\min_u\max_v$ 换成 $\max_v\min_u$ 的对应方程。）
\end{theorem}

把 \cref{eq:dg-hji} 与前置桥接的 HJB 方程 \cref{eq:dg-hjb} 并排看，结构完全平行——唯一区别是单人的 $\min_u$ 变成了博弈的 $\min_u\max_v$。

\begin{note}
\textbf{多视角理解（HJB $\leftrightarrow$ HJI）.}\;HJB 与 HJI 像同一台机器换了个旋钮。HJB 在每个状态点问“我这一步该怎么走，使（当前代价 + 未来代价变化）最小”；HJI 问“我求最小、对手求最大，在这个对抗下这一步的鞍点值是多少”。相似之处：都是值函数对状态梯度 $\nabla_x V$（协态）与动力学 $f$ 做内积、再叠加瞬时代价 $L$，构成哈密顿量。不同之处：HJI 的哈密顿量内部是一个 $\min$-$\max$ 博弈而非单个 $\min$，这一改动把“最优”变成了“鞍点”，也带来了 Isaacs 条件这个全新问题。类比的边界：不要因此以为 HJI“只是 HJB 多套一层优化”——多出来的这层让解的存在唯一性、可微性都变复杂了（\cref{sec:dg-viscosity}）。
\end{note}

\paragraph{哈密顿量与最优策略。}
定义哈密顿量
\begin{equation}
  H(x,\lambda)=\min_u\max_v\big[\,\lambda\cdot f(x,u,v)+L(x,u,v)\,\big],\qquad \lambda=\nabla_x V.
  \label{eq:dg-ham}
\end{equation}
给定值函数，双方的瞬时最优策略由这个内部博弈的鞍点给出：$u^*,v^*=\arg\min_u\max_v[\lambda\cdot f+L]$。也就是说，求解 HJI 方程同时给出了值函数和最优反馈策略——和单人 HJB 一样，这是动态规划方法相对 PMP 的核心优点：PMP 给的是单条最优轨迹（开环），HJI 给的是全状态空间的反馈律 $u^*(x,t)$。反馈律的好处在 \cref{sec:dg-reach} 会再强调：有扰动时反馈策略能“见状态调整”，比开环鲁棒得多。

\begin{note}
\textbf{新手常见困惑（Q\&A）.}\;\textbf{Q1：为什么 \cref{eq:dg-dpp} 里那个 $\inf_\gamma\sup_u$（对整段策略求极值）到极限就变成了逐点的 $\min_u\max_v$（只对此刻动作求极值）？}\;因为取了 $\delta\to0$。$\inf_\gamma\sup_u$ 是对 $[t,t+\delta]$ 这一小段上的策略函数求极值；当 $\delta\to0$，这一段塌缩成一个时刻，“策略函数”退化为“此刻的一个动作 $(u,v)$”，于是对函数的 $\inf\sup$ 就变成了对当前动作的 $\min\max$。未来的全部影响已被 $\nabla_x V\cdot f$ 那一项（值函数的梯度）打包带走。\textbf{Q2：为什么 HJI 是终端条件、向“过去”解？}\;因为代价定义在未来：$V(x,t)$ 是“从 $(x,t)$ 到终点 $T$ 的最优剩余代价”，终点处没有剩余博弈，值就等于终端代价 $g$——这是唯一能直接写出的边界，其余时刻靠它向回递推。\textbf{Q3：HJI 给反馈、PMP 给开环，为什么？}\;PMP 是\emph{必要条件}，沿\emph{一条}最优轨迹给出协态方程和极值条件，解出开环轨迹 $u^*(t)$；HJI 是\emph{充分}的动态规划方程，在\emph{整个状态空间}解出值函数、进而给反馈律 $u^*(x,t)$。开环被扰偏了就失效，反馈律能实时纠正。代价是 HJI 要解全空间的 PDE（维度诅咒，\cref{sec:dg-curse}）。\textbf{Q4：HJI 的协态 $\lambda=\nabla_x V$ 和 PMP 的协态是一回事吗？}\;本质同一个：PMP 的协态 $\lambda(t)$ 沿最优轨迹恰好等于 $\nabla_x V(x^*(t),t)$，都是“状态的影子价格”；HJI 是这个关系在全空间的版本，PMP 是它沿单条轨迹的切片。
\end{note}

\subsection{Isaacs 条件：博弈到底有没有“值”}
\label{sec:dg-isaacs-cond}
这是本节的关键。下值 $V^-$ 对应逐点哈密顿量的 $\max_v\min_u$（信息优势给最小化方 $u$），上值 $V^+$ 对应 $\min_u\max_v$（信息优势给最大化方 $v$）；\cref{prop:dg-maxmin} 已证 $\max_v\min_u\le\min_u\max_v$，即 $V^-\le V^+$。\textbf{当两者相等时，博弈有明确定义的“值”，且存在鞍点策略}——这就是 \textbf{Isaacs 条件（鞍点条件）}：
\begin{equation}
  \min_u\max_v\big[\,\lambda\cdot f(x,u,v)+L\,\big]=\max_v\min_u\big[\,\lambda\cdot f(x,u,v)+L\,\big].
  \label{eq:dg-isaacs-cond}
\end{equation}

\begin{insight}
Isaacs 条件成立，意味着“谁先暴露策略”不影响结果——上值等于下值，博弈有唯一的值，先后手优势消失。它在博弈论里对应 von Neumann 的 minimax 定理在逐点哈密顿量层面的体现：内部那个关于 $(u,v)$ 的静态博弈存在鞍点。
\end{insight}

\paragraph{什么时候成立？}
一个充分（非必要）条件是哈密顿量内部对 $u,v$ \textbf{可分离}。如果动力学与代价能写成
\begin{equation}
  \lambda\cdot f+L=\underbrace{[\,\text{只含 }u\,]}_{\text{P1 管}}+\underbrace{[\,\text{只含 }v\,]}_{\text{P2 管}}+[\,\text{不含 }u,v\,],
  \label{eq:dg-separable}
\end{equation}
那么 $\min_u$ 只作用在“只含 $u$ 的项”、$\max_v$ 只作用在“只含 $v$ 的项”，二者作用于互不重叠的部分、互不干涉；既然先对哪一部分取极值都不影响另一部分，交换 $\min_u$ 与 $\max_v$ 的顺序就不改变结果，Isaacs 条件自动成立。大量追逃博弈恰好满足这个可分离性，因为 $u$ 和 $v$ 常常各自\emph{线性地}出现在动力学里（如 $f=a(x)+b(x)u+c(x)v$），且代价对二者解耦。更一般地，只要内部博弈对 $u$ 凸、对 $v$ 凹（鞍点存在的经典充分条件），Isaacs 条件就成立。失效的情形：当 $u$ 与 $v$ 在哈密顿量里\emph{乘在一起}（如出现 $u^\top \mathbf M v$ 这样的双线性耦合且 $\mathbf M$ 不定），内部博弈未必有鞍点，$\min\max>\max\min$ 严格成立，博弈只有上下值而没有统一的“值”；此时强行套 \cref{eq:dg-hji} 会得到一个并不对应任何一方真实最优的“伪值”。

\begin{example}[一个算到底的一维追逃]
\label{ex:dg-1d}
设相对位置 $x\in\mathbb{R}$，动力学 $\dot x=-u+v$，$u\in[-1,1]$（追捕者）、$v\in[-1,1]$（逃跑者），无运行代价 $L=0$，终端代价 $g(x)=|x|$。协态 $\lambda=\partial V/\partial x$，哈密顿量
\[
  H(x,\lambda)=\min_{u\in[-1,1]}\max_{v\in[-1,1]}\ \lambda(-u+v)=\min_u[-\lambda u]+\max_v[\lambda v].
\]
$u$ 项与 $v$ 项已分离，故 $\min\max=\max\min$，\textbf{Isaacs 条件成立}。逐项求极值：
\[
  \min_{u}(-\lambda u)=-|\lambda|\Rightarrow u^*=\operatorname{sign}(\lambda),\qquad \max_{v}(\lambda v)=|\lambda|\Rightarrow v^*=\operatorname{sign}(\lambda).
\]
于是 $H=-|\lambda|+|\lambda|=0$。代入 HJI 方程 $-\partial_t V=H=0$，得 $\partial_t V=0$——值函数不随时间变化，恒等于终端代价 $V(x,t)=|x|$。物理意义清晰：双方“操控权”在动力学里完全对称（$-u+v$，系数大小相同），全力对抗的净效果相互抵消，相对距离保持不变。最优策略 $u^*=v^*=\operatorname{sign}(x)$ 也直白：$x>0$ 时追捕者取 $u=+1$ 想把 $x$ 拉小、逃跑者取 $v=+1$ 想把 $x$ 推大，两者打平。注意 $V=|x|$ 在 $x=0$ 处不可微——这正是 \cref{sec:dg-viscosity} 要处理的“棱”，也提醒我们 $\lambda=\operatorname{sign}(x)$ 在原点没有定义、需要黏性解框架来收尾。
\end{example}

\begin{pitfall}[思维陷阱：默认 $\min_u\max_v=\max_v\min_u$]
写 Isaacs 方程时随手把 $\min\max$ 写成 $\max\min$，或反过来，以为两者本来就一样。\textbf{后果}：在 Isaacs 条件不成立的博弈里，你算出的值函数既不是上值也不是下值，得到的“最优策略”对任何一方都不是真的最优；仿真里表现为双方行为诡异、解对初始化异常敏感。\textbf{根因}：$\max\min\le\min\max$ 是普适不等式（\cref{prop:dg-maxmin}），等号只在 Isaacs 条件（鞍点存在）下成立。\textbf{正解}：写方程前先验证 Isaacs 条件——检查哈密顿量对 $u,v$ 是否可分离（或凸-凹）。可分离时随便用哪个顺序都行；不可分离时必须明确你算的是上值还是下值，并据此选 $\min\max$ 或 $\max\min$。
\end{pitfall}

\begin{pitfall}[概念误区：忽略信息结构，以为只有“一个”博弈值]
把零和博弈当成“双方各求一次极值”，不区分谁能看到谁的当前动作。\textbf{后果}：在需要鲁棒安全保证的场景（\cref{sec:dg-reach}），用错信息结构会高估安全集——让对手（干扰）反而吃亏，得到偏乐观、不安全的结果。\textbf{根因}：博弈的值依赖信息结构（开环/反馈/非预期策略）；安全验证里要给对手瞬时信息优势（非预期策略），才能保证“最坏情况”真的被覆盖。\textbf{正解}：明确写出谁用非预期策略、谁先承诺。安全验证默认让干扰拥有瞬时信息优势（取下值），这样得到的安全集是保守、可信的。检验法：问“如果对手能多看我一眼，结果会不会更糟？”——会，就说明你的信息结构给安全留了漏洞。
\end{pitfall}

\begin{practice}
\begin{enumerate}
  \item （手推，照着 \cref{ex:dg-1d} 做一遍）把一维追逃改成“操控权不对称”：$\dot x=-2u+v$（追捕者更强），其余不变。重新写哈密顿量，求 $u^*,v^*$ 与 $H(x,\lambda)$，解出 $V(x,t)$，并解释结果（追捕者能不能稳定缩小距离？）。
  \item （判定 Isaacs 条件）对 $H(u,v)=u^2-v^2+\lambda(u+v)$（$u,v\in\mathbb{R}$）与 $H(u,v)=uv+\lambda u$（$u,v\in[-1,1]$），分别判断 $\min_u\max_v$ 是否等于 $\max_v\min_u$，给出鞍点（若存在）或说明为何不存在。把结论对应到“可分离 / 凸-凹 / 双线性耦合”哪种情形。
  \item （概念联系）解释为什么“安全验证”天然是零和博弈：谁是 $\min$ 方、谁是 $\max$ 方？为什么给干扰非预期策略对应“最坏情况安全”？再用 \cref{eq:dg-maxmin} 说明：若错误地让控制占信息优势，安全集会被高估还是低估？
\end{enumerate}
\end{practice}

% ════════════════════════════════════════════════════════════════
\section{HJI 偏微分方程与黏性解：值函数为什么不可微}
\label{sec:dg-viscosity}

\paragraph{动机：从静态方程到可数值求解的演化方程。}
\cref{sec:dg-isaacs} 的 \cref{eq:dg-hji} 是一个终端值偏微分方程：给定终端时刻 $V(x,T)=g(x)$，向时间负方向推出每个时刻的值函数。直觉上只要从终端往回积分这个 PDE 就能得到 $V$。但这里藏着一个数学陷阱：\textbf{这个积分过程会产生不可微的“棱”}。而一旦 $\nabla_x V$ 在某处不存在，方程里的 $\nabla_x V\cdot f$ 就失去了意义，经典意义下的“解”根本不存在。\cref{ex:dg-1d} 那个一维追逃已经撞上这件事：值函数 $V(x,t)=|x|$ 在 $x=0$ 处有一道棱，协态 $\lambda=\operatorname{sign}(x)$ 在原点没定义。

\paragraph{反面：为什么值函数会长出“棱”。}
在追逃博弈里，从某个初始状态出发，可能有两条不同的最优策略给出相同的代价（比如向左绕和向右绕都一样快地被追上）。在这两类初始状态的分界面上，值函数 $V$ 连续，但它的梯度从一侧跳到另一侧——出现一道折痕，就像两个斜面相交的屋脊。

\begin{note}
\textbf{多视角理解（值函数的“棱” $\leftrightarrow$ 距离场的脊线）.}\;这道棱不是病态特例，而是普遍现象。一个直观类比：考虑平面上一个非凸障碍物的“到障碍物的距离场”——在障碍物凹进去的地方背后，会形成一条脊线，脊线上每一点到障碍边界有两个等距最近点，距离函数在那里有折痕、不可微。HJI 值函数的棱与此同源：它本质也是一种“代价场”，在“有多条等价最优路径”的地方就会起脊。相似之处：二者都是“多个等距/等价竞争者在某处打平”导致的不可微。不同之处：距离场只关于几何，而 HJI 值函数还编码了对抗动力学。这个类比提醒我们：要求 HJI 的解处处可微，就像要求一个非凸物体的距离场处处可微一样不现实。
\end{note}

\subsection{一个能算到底的反例：eikonal 方程}
\label{sec:dg-eikonal}
抛开时间，看最简单的一阶 HJ 方程——一维 eikonal 方程
\begin{equation}
  |u'(x)|=1\quad\text{在 }(-1,1)\text{ 上},\qquad u(-1)=u(1)=0.
  \label{eq:dg-eikonal}
\end{equation}
它的物理意义是“到边界 $\{-1,1\}$ 的距离场应满足梯度模长为 1”。\textbf{问题在于：几乎处处满足 $|u'|=1$ 的连续函数有无穷多个}。最简单的两个：$u_+(x)=1-|x|$（一个尖朝上的帐篷，在 $x=0$ 达到峰值 1，两臂斜率 $\mp1$）；$u_-(x)=|x|-1$（一个尖朝下的山谷，在 $x=0$ 达到谷底 $-1$）。除此之外还能拼出各种锯齿状的解（在 $(-1,1)$ 里来回折、每段斜率 $\pm1$）。单看方程它们平起平坐。但\textbf{物理上正确的只有一个}：$u_+(x)=1-|x|$，它恰好是到边界的距离，处处非负、在中点最大。$u_-$（负的距离）和那些锯齿解都是数学假象。怎么用一个原则把 $u_+$ 挑出来、把其余否决掉？这就是黏性解。

\subsection{循序渐进：为什么朴素的弱解都行不通}
\label{sec:dg-weak-fail}
在跳到黏性解的定义前，先放慢脚步，把“我们到底需要一个什么样的解”想清楚。面对一个解会起棱的方程，最自然的做法是依次试几种现成的“解”概念——看着它们一个个失败，黏性解为什么非它不可就水落石出了。

\textbf{尝试一：经典解（处处可微）。}\;最理想的当然是要求解处处一阶连续可微、逐点满足方程。但值函数普遍有不可微的棱（eikonal 的 $1-|x|$ 在 $x=0$、追逃值函数在散射面上）。棱处 $\nabla V$ 不存在，方程里的 $\nabla V\cdot f$ 无从谈起——\textbf{经典解根本不存在}。

\textbf{尝试二：分布意义的弱解（像线性 PDE 那样）。}\;线性 PDE 里有个标准招数：把方程乘上一个光滑测试函数、做分部积分，把导数“转移”到测试函数身上，于是即便解本身不可微也能定义弱解（分布解）。这招对线性方程无往不利。可 HJ/HJI 方程对梯度是\textbf{非线性}的——哈密顿量里是 $|\nabla V|$、$H(\nabla V)$ 这种非线性项；分部积分搬不动一个非线性的梯度函数，分布之间也不能做非线性的复合或相乘。\textbf{分布弱解这套机器对非线性 HJ 根本启动不了。}

\textbf{尝试三：“几乎处处”满足方程。}\;那退一步：棱是测度零的集合，要求解在可微的地方（也就是几乎处处）逐点满足方程，行不行？这个概念是良定义的——但 eikonal 反例已经演示过它的致命伤：\textbf{a.e. 解不唯一}（$1-|x|$、$|x|-1$、各种锯齿解都几乎处处满足 $|u'|=1$）。方程加“a.e.”这个条件太弱，选不出物理正确的那一个。

三次尝试、三种失败，恰好框定了我们需要的解必须同时满足三个看似矛盾的要求（\cref{tab:dg-weak}）。\textbf{黏性解正是同时满足这三条的那个概念}：它弱到允许棱、绕开了分布对非线性的无能、又强到能唯一确定物理解。

\begin{table}[H]\centering\small
\caption{各种“解”概念的能力对照（$\checkmark$ 满足 / $\times$ 不满足）}
\label{tab:dg-weak}
\begin{tabular}{lcccc}
\toprule
要求 & 经典解 & 分布弱解 & a.e. 解 & 我们想要的 \\
\midrule
允许棱（不要求处处可微） & $\times$ & $\checkmark$ & $\checkmark$ & $\checkmark$ \\
适用于对梯度非线性的方程 & $\checkmark$ & $\times$ & $\checkmark$ & $\checkmark$ \\
唯一（能选出物理解） & $\checkmark$ & --- & $\times$ & $\checkmark$ \\
\bottomrule
\end{tabular}
\end{table}

\subsection{慢动作：怎样在不可微的点上“满足”一个方程}
\label{sec:dg-slow}
黏性解背后那个真问题，对没学过偏微分方程的工程师是全章最反直觉的一处，值得多花几分钟。\textbf{问题摆明白}：方程里有 $\nabla_x V$，在值函数的棱上 $\nabla_x V$ 根本不存在，那么“$V$ 在棱上满足方程”到底是什么意思？\textbf{关键的转念}：你不能对 $V$ 求导，但你可以拿一根\textbf{光滑}的函数 $\phi$ 去贴住 $V$，然后对 $\phi$ 求导——$\phi$ 是你自己挑的光滑函数，它的导数处处存在。于是“方程在 $x_0$ 处成立”这件对 $V$ 做不到的事，被转嫁给了在 $x_0$ 处贴住 $V$ 的那根光滑 $\phi$。这就是黏性解的全部诡计。

\textbf{慢动作之一：从上方贴。}\;想象在棱点 $x_0$ 处，一根光滑、开口向下的“碗” $\phi$ 从\emph{上方}压下来，一路下降到刚好碰到 $V$（即 $\phi\ge V$ 在 $x_0$ 附近、$\phi(x_0)=V(x_0)$；等价地 $V-\phi$ 在 $x_0$ 取局部极大）。在这个碰触点，$\phi$ 有一个确定的斜率 $\nabla\phi(x_0)$。黏性解要求方程对这个斜率成立\emph{一个方向}的不等式（次解条件）。\textbf{为什么是“一个方向”、不是等号？}\;因为从上方贴的 $\phi$，斜率被 $V$ 的形状“卡”在一个范围里，不是唯一值：在一个向上的尖峰处，能从上方贴住的光滑 $\phi$ 的斜率可以是两条臂斜率（$+1$ 与 $-1$）之间的任意值。对所有这些 $\phi$ 都要求方程取等号太强、根本做不到；只要求它们都满足同一方向的不等式，就刚刚好。

\textbf{慢动作之二：从下方贴。}\;对称地，一根开口向上的光滑“碗” $\phi$ 从\emph{下方}顶上来碰到 $V$（$\phi\le V$ 附近、$\phi(x_0)=V(x_0)$，即 $V-\phi$ 取局部极小），给出\emph{另一个方向}的不等式（超解条件）。两面一卡，棱就定住了：一个点若从上方贴得一个方向、从下方贴得另一个方向、两者都满足，它就是黏性解。打个比方：像用两块光滑模板从上下两面卡一个有棱的零件——单独一面卡不死棱的位置，两面同时卡，棱就被唯一确定了。

\begin{definition}[黏性解]
\label{def:dg-viscosity}
把方程统一写成 $F(x,\nabla u)=0$（eikonal 即 $F=|\nabla u|-1$；HJI 即 $F=\partial_t V+H$）。
\begin{itemize}
  \item \textbf{黏性次解（subsolution）}：对任意光滑 $\phi$，只要 $u-\phi$ 在某点 $x_0$ 取\emph{局部极大}（即 $\phi$ 从\emph{上方}触碰 $u$），就要求 $F(x_0,\nabla\phi(x_0))\le 0$；
  \item \textbf{黏性超解（supersolution）}：对任意光滑 $\phi$，只要 $u-\phi$ 在某点 $x_0$ 取\emph{局部极小}（即 $\phi$ 从\emph{下方}触碰 $u$），就要求 $F(x_0,\nabla\phi(x_0))\ge 0$；
  \item \textbf{黏性解}：同时是次解和超解。
\end{itemize}
在 $u$ 可微的点，上下两个条件合起来退化为 $F(x,\nabla u)=0$（经典方程）；在棱上，则由上下夹逼挑出唯一正确的解。
\end{definition}

黏性解理论由 Crandall 与 Lions（1983）为 Hamilton--Jacobi 方程奠定\cite{crandall1983viscosity}（Lions 后来因相关工作获菲尔兹奖）。

\begin{derivation}[两个不等号方向来自“消失黏性”]
\label{deriv:dg-vanishing}
给方程加一个小扩散项，考虑 $F(x,\nabla u^\varepsilon)=\varepsilon\,\Delta u^\varepsilon$，其光滑解为 $u^\varepsilon$。设 $\phi$ 光滑且 $u-\phi$ 在 $x_0$ 取严格局部极大；当 $\varepsilon$ 小时，$u^\varepsilon-\phi$ 在附近某点 $x_\varepsilon\to x_0$ 也取局部极大。在极大点处，一阶条件给 $\nabla u^\varepsilon=\nabla\phi$，二阶条件给 $\Delta(u^\varepsilon-\phi)\le 0$，即 $\Delta u^\varepsilon\le\Delta\phi$。代入正则化方程：
\[
  F(x_\varepsilon,\nabla\phi)=\varepsilon\,\Delta u^\varepsilon\le\varepsilon\,\Delta\phi\ \xrightarrow{\ \varepsilon\to 0\ }\ F(x_0,\nabla\phi)\le 0.
\]
这就推出了次解的 $\le 0$。超解（从下方触碰、局部极小）对称地给出 $\ge 0$。所以“从上方碰得 $\le 0$、从下方碰得 $\ge 0$”是消失黏性极限的逻辑后果，不是人为规定。
\end{derivation}

\begin{insight}
黏性解不是“放松要求凑合用”的妥协，而是抓住了正确解的本质特征。“黏性”二字正来自上面这个极限：给方程加一个无穷小的黏性（扩散）项，扩散会把棱磨光、得到光滑解 $u^\varepsilon$；让 $\varepsilon\to0$，$u^\varepsilon$ 收敛到的那个极限，恰好就是黏性解。所以黏性解 = “无穷小黏性正则化的极限”——它是物理上、数值上都站得住的那一个，也正是带数值耗散的单调格式（\cref{sec:dg-numerics}）会收敛到的解。
\end{insight}

\subsection{用黏性解否决假解：回到 eikonal}
\label{sec:dg-eikonal-judge}
现在用这套机制审判前面那两个候选解。关键都发生在尖点 $x=0$（别处都光滑、自动满足方程）。

\textbf{审判 $u_+(x)=1-|x|$（尖朝上的峰）。}\;能否有光滑 $\phi$ 从\emph{下方}触碰（$u_+-\phi$ 局部极小）？要 $\phi\le u_+$ 附近、$\phi(0)=u_+(0)$。但 $u_+$ 在 $0$ 处是个向上的尖峰，两臂以斜率 $\mp1$ 下降；一个光滑函数要从下方贴住这个尖峰，需同时满足 $x>0$ 侧斜率 $\le-1$、$x<0$ 侧斜率 $\ge+1$，自相矛盾——\textbf{没有光滑 $\phi$ 能从下方触碰}。于是超解条件在 $x=0$ 处\emph{空洞地满足}（没有需要检验的 $\phi$）。而从上方触碰的 $\phi$，其斜率 $\phi'(0)\in[-1,1]$，次解条件 $|\phi'(0)|-1\le 0$ 成立。所以 $u_+$ 同时是次解和超解，\textbf{是黏性解}。

\textbf{审判 $u_-(x)=|x|-1$（尖朝下的谷）。}\;现在 $u_-$ 在 $0$ 处是向下的尖谷。取一条平缓的光滑曲线从\emph{下方}触碰谷底，$\phi'(0)=0$ 是允许的（它确实能贴在谷底下方、在 $0$ 处相切）。此时 $u_--\phi$ 在 $0$ 取局部极小，触发超解条件：要求 $F(0,\phi'(0))=|\phi'(0)|-1\ge 0$。但 $|0|-1=-1<0$，\textbf{违反超解条件}。所以 $u_-$ 不是黏性超解，\textbf{不是黏性解、被否决}。

\begin{note}
\textbf{对比性思维（峰 vs 谷的命运）.}\;同样是一道尖，朝上的峰能当黏性解、朝下的谷不能，差别全在“哪个方向能被光滑函数触碰”。向上的尖峰挡住了来自下方的所有光滑触碰（超解条件空洞满足）；向下的尖谷却任由光滑函数从下方贴入，于是被超解条件抓个正着。这背后是一条几何直觉：黏性解只允许“凸向正确方向”的棱。对距离场而言，正确的棱是“远离边界的脊线”（向上的峰），所以 $u_+=1-|x|$（距离）胜出。那些锯齿解里也藏着向下的谷，会被同样的机制逐一否决——这就是黏性解把无穷多个 a.e. 解收缩到唯一一个的方式。
\end{note}

\begin{note}
\textbf{多视角理解（HJ 的棱 $\leftrightarrow$ 守恒律的激波）.}\;黏性解的概念在另一个领域有个孪生兄弟。把一维 HJ 方程 $u_t+H(u_x)=0$ 对 $x$ 求导、令 $w=u_x$，就得到一个\emph{标量守恒律} $w_t+H(w)_x=0$。在这个映射下：HJ 解 $u$ 的不可微“棱”（$u_x$ 跳变处）恰好对应守恒律解 $w$ 的“激波”（间断）；HJ 的黏性解条件恰好对应守恒律的\emph{熵条件}（哪些激波物理可容许）。相似之处：两者都面对“a.e. 弱解不唯一”的困境，都用一个额外条件（黏性/熵，都来自加无穷小扩散再取极限）挑出唯一物理解。不同之处：这个精确对应只在一维成立；高维里 HJ 与守恒律的关系松散得多，不能逐字搬用。但它给了一个有力直觉：HJI 值函数的棱，就是“代价场里的激波”。
\end{note}

\subsection{值函数 = HJI 的唯一黏性解}
\label{sec:dg-unique}
把这一切钉死的关键定理来自 Evans 与 Souganidis（1984）\cite{evans1984differential}。

\begin{theorem}[博弈值 = 唯一黏性解（Evans--Souganidis 1984）]
\label{thm:dg-evans}
在 Isaacs 条件下，零和微分博弈的上值与下值分别是相应 HJI 方程的\textbf{唯一黏性解}，并有相应的表示公式。唯一性的机制是\textbf{比较原理（comparison principle）}：若 $u$ 是次解、$w$ 是超解，且在边界/终端上 $u\le w$，则在整个区域上 $u\le w$。
\end{theorem}
\begin{proof}[唯一性（比较原理夹逼）]
设有两个黏性解 $V_1,V_2$，终端条件相同（都等于 $g$）。把 $V_1$ 当次解、$V_2$ 当超解用比较原理 $\Rightarrow V_1\le V_2$；反过来把 $V_2$ 当次解、$V_1$ 当超解 $\Rightarrow V_2\le V_1$。两者一夹，$V_1=V_2$。次解/超解那两个看似古怪的单边条件，正是为了让这个夹逼论证能跑通而设计的。
\end{proof}

这个结果的意义是双重的：(1) 它保证 \cref{sec:dg-isaacs} 写出的 HJI 方程\textbf{在黏性解意义下有且仅有一个解}，方程不再因不可微和弱解不唯一而失效；(2) 它把“博弈的值”（一个由策略定义的对象）与“HJI 方程的黏性解”（一个由 PDE 定义的对象）严格等同起来——这正是后续一切数值方法的合法性根基：在网格上算 PDE 的黏性解，就是在算博弈的值。到这里，从动机到严格基础的链条闭合了：\cref{sec:dg-isaacs} 用动态规划把博弈写成 HJI 方程，\cref{sec:dg-viscosity} 用黏性解理论保证这个方程有且仅有一个物理正确的解。剩下的问题是——怎么把它算出来？这就引出 \cref{sec:dg-reach} 的可达性表述与 \cref{sec:dg-numerics} 的数值方法。特别地，\cref{sec:dg-numerics} 的数值格式必须“尊重黏性”：用带数值耗散的单调格式（如 Lax--Friedrichs 型、WENO），才能在棱处稳定收敛到黏性解，而不是收敛到某个被否决的假解。

\begin{note}
\textbf{新手常见困惑（Q\&A）.}\;\textbf{Q1：为什么偏偏叫“黏性”解？}\;名字直接来自消失黏性极限（\cref{deriv:dg-vanishing}）。加的 $\varepsilon\Delta V$ 是一个扩散/黏性项（就像流体里的黏性会把激波抹平），它把值函数的棱磨光成光滑的 $V_\varepsilon$；让 $\varepsilon\to0$，得到的极限就是黏性解。\textbf{Q2：加的那个扩散项可以随便选吗？}\;关键的好性质正在这里：对一大类正则化，极限都收敛到同一个黏性解。正是这种“对正则化方式不敏感”的稳健性，让黏性解成为良定义的概念。数值上对应的事实是：不同的单调（带耗散）格式都收敛到同一个黏性解。\textbf{Q3：博弈的值函数一定是黏性解吗？}\;一定是——这正是 \cref{thm:dg-evans} 的内容。在 Isaacs 条件下，博弈的值\emph{恰好就是} HJI 方程的唯一黏性解。这座桥是双向的：要算博弈的值就去求 HJI 的黏性解，求得了黏性解就得到了博弈的值。\textbf{Q4：比较原理凭什么保证唯一？}\;直觉是一个“夹逼”，见上面证明。
\end{note}

\begin{pitfall}[概念误区：默认 HJI 值函数处处可微，直接用 $\nabla V$]
把值函数当成光滑函数，在推导或写数值格式时随手对 $V$ 求梯度、求二阶导。\textbf{后果}：在棱附近梯度不存在，数值上 $\nabla V$ 剧烈震荡甚至发散；用中心差分等不带耗散的格式会在折痕处产生伪振荡，可达集边界长出毛刺或断裂。\textbf{根因}：HJI 值函数普遍含不可微的棱（多条等价最优路径的分界面），它是黏性解而非经典解。\textbf{正解}：用为黏性解设计的、带数值耗散的单调格式（\cref{sec:dg-numerics} 的 WENO + 合适的数值哈密顿量），它们在折痕处稳定收敛到正确的黏性解；理论分析时引用黏性解框架而非假设可微。
\end{pitfall}

\begin{practice}
\begin{enumerate}
  \item （亲手否决一个锯齿解）在 $(-1,1)$ 上构造 eikonal $|u'|=1$、$u(\pm1)=0$ 的一个锯齿解（比如在 $[-1,0]$ 升、$[0,1]$ 降之外再多折一次）。指出它在哪个尖点会被黏性解的哪个条件（次解还是超解）否决，照着 $u_-$ 的审判过程写一遍。
  \item （构造直觉）取 \cref{ex:dg-1d} 的值函数 $V(x,t)=|x|$。说明它在 $x=0$ 处为何不可微、对应“哪条最优路径在此打平”，并验证它是相应 HJI 方程的黏性解（提示：原点是向上还是向下的尖？）。
  \item （文献定位）查阅 Evans--Souganidis（1984）的主结论陈述，用 2--3 句话写出：他们证明了博弈的上/下值与 HJI 方程的什么解等同？比较原理在唯一性证明中起什么作用？这个等同为什么是后续数值方法的合法性基础？
\end{enumerate}
\end{practice}

% ════════════════════════════════════════════════════════════════
\section{可达性分析：BRS、BRT 与 reach--avoid}
\label{sec:dg-reach}

\paragraph{动机：把“安全吗”变成“从哪些状态出发会出事”。}
安全验证要回答的是：从当前状态出发，在最坏情况干扰下，系统会不会进入危险区？把这个问题反过来问更好算——\textbf{有哪些初始状态，会不可避免地（在对手最优、我方也最优的对抗下）落入危险集？}这个“会出事的初始状态集合”就是可达性分析的产物。它一旦算出来，安全控制就有了依据：只要系统不在这个集合里，就有控制能保证安全；它的补集就是安全集。注意这里用的正是 \cref{sec:dg-why} 的“零和 $\leftrightarrow$ 鲁棒控制”视角：把最坏情况干扰当成对抗玩家。

\paragraph{历史与背景。}
把可达集计算与 HJI 方程严格挂钩、并给出网格数值方法的，是 Mitchell、Bayen 与 Tomlin（2005）——“连续动态博弈的可达集的时变 Hamilton--Jacobi 表述”，这是可达性领域的奠基性工程文献\cite{mitchell2005time}。本节的公式与角色分析主要依据这条线后来的标准教程整理（Bansal、Chen、Herbert、Tomlin，2017，CDC）\cite{bansal2017hamilton}，它把 BRS/BRT 与角色取法讲得最清楚。“既要到达目标、又要全程避开危险”的 reach--avoid 表述由 Margellos 与 Lygeros（2011）系统化\cite{margellos2011hamilton}，Fisac 等（2015）进一步推广到时变的动力学、目标与约束\cite{fisac2015reach}。

\paragraph{理论：把布尔问题编码成连续值函数。}
可达性本质是个\emph{布尔问题}（系统到没到危险集），布尔问题不好做连续优化——这正是 \cref{sec:dg-why} 末尾提到的“类型博弈 vs 程度博弈”。水平集方法的核心技巧（已在前置桥接铺垫）是把类型博弈化成程度博弈：用一个 Lipschitz 函数 $g(x)$ 把目标集 $\mathcal{T}$ 编码成它的零下水平集，
\begin{equation}
  x\in\mathcal{T}\ \Longleftrightarrow\ g(x)\le 0,\qquad g=\text{到 }\mathcal{T}\text{ 的带符号距离},
  \label{eq:dg-sdf}
\end{equation}
然后把代价取成纯终端代价 $J=g(x(0))$（无运行代价）。这样“到没到目标”就等价于“终端代价是否 $\le0$”，布尔问题被一个连续值函数接管。下面采用可达性文献的标准约定：时间反向，$t<0$ 表示“距终端还有 $|t|$ 时长”，系统在 $[t,0]$ 上演化。

\subsection{后向可达集（BRS）：恰好在终点到达}
\label{sec:dg-brs}
设玩家 1（控制 $a$）想\emph{远离}目标、玩家 2（干扰 $b$）想\emph{驱入}目标。BRS 是这样一组初始状态：从中出发，干扰能（用非预期策略）在时刻 $0$ 把系统送进目标，不管控制怎么努力：
\begin{equation}
  \mathcal{R}(t)=\Big\{x:\ \exists\,\gamma,\ \forall a(\cdot),\ \zeta\big(0;x,t,a(\cdot),\gamma[a](\cdot)\big)\in\mathcal{T}\Big\}.
  \label{eq:dg-brs-def}
\end{equation}
对应的值函数 $V(x,t)$（取下值，给干扰瞬时信息优势）是下面 HJI 方程的黏性解：
\begin{align}
  D_t V(x,t)+H\big(x,\nabla V\big)&=0,\qquad V(x,0)=g(x), \label{eq:dg-brs-pde}\\
  H(x,\lambda)&=\max_{a}\ \min_{b}\ \lambda\cdot f(x,a,b). \label{eq:dg-brs-ham}
\end{align}
于是 BRS 就是值函数的零下水平集：$\mathcal{R}(t)=\{x:V(x,t)\le 0\}$。注意 \cref{eq:dg-brs-pde} \textbf{不含任何额外的冻结项}——这是“集”的方程。

\paragraph{哈密顿量里 max/min 的取法：一句口诀。}
\cref{eq:dg-brs-ham} 里 $\max_a\min_b$ 的取法不是随手定的，由角色决定，这一点必须讲透。规则极简，记住它能省掉无数次符号纠结。

\begin{note}
\textbf{系统性分类（哈密顿量里 max/min 的取法）.}\;看你对每个玩家问的是“存在”还是“任意”：
\begin{itemize}
  \item 求一个控制的\textbf{存在性}（$\exists a$，“只要有一个动作能……”）$\to$ 在哈密顿量里对它取 $\min$；
  \item 集合要对该控制的\textbf{所有取值}都成立（$\forall a$，“不管它怎么动……”）$\to$ 对它取 $\max$。
\end{itemize}
在 BRS 定义里，我们要“$\exists$ 干扰 $b$”（干扰能驱入）且“$\forall$ 控制 $a$”（不管控制怎么挡），所以哈密顿量是 $\max_a\min_b$，与 \cref{eq:dg-brs-ham} 一致。换一种角色——“$\exists$ 控制把系统送进好的目标，$\forall$ 干扰捣乱”——就翻成 $\min_a\max_b$。这套“存在取 min、任意取 max”的口诀对 reach、avoid、reach--avoid 一律适用。
\end{note}

\begin{example}[单积分器的可达管，含解析验证]
\label{ex:dg-integrator}
抛开干扰先看最干净的情形。设一维单积分器 $\dot x=u$，控制 $u\in[-1,1]$（最大速度 1），目标 $\mathcal{T}=\{|x|\le1\}$，带符号距离 $g(x)=|x|-1$。控制想\emph{到达}目标，故“$\exists u$”、哈密顿量取 $\min$：$H(x,\lambda)=\min_{u\in[-1,1]}\lambda u=-|\lambda|$。直觉上，在时间 $|t|$ 内最大速度 1 能覆盖距离 $|t|$，所以能到达目标的初始状态就是“离目标不超过 $|t|$”的那些，即 $\{|x|\le1+|t|\}$，对应值函数 $V(x,t)=|x|-1-|t|$。验证（用下面 BRT 的冻结项方程，因为“在窗内到达”是管）：$\partial_x V=\operatorname{sign}(x)$ 故 $|\partial_x V|=1$、$H=-1$；对 $t<0$ 有 $\partial_t V=+1$；于是 $D_t V+\min[0,H]=1+\min[0,-1]=0$，且 $V(x,0)=|x|-1=g(x)$，正确。可达管 $\mathcal{R}_{\text{tube}}(t)=\{|x|\le1+|t|\}$ 随 $|t|$ 线性扩张，完全符合“速度 × 时间”的直觉。

加上对抗干扰会怎样？设 $\dot x=u+d$，控制 $u\in[-1,1]$ 想远离\emph{危险}目标 $\{|x|\le1\}$、干扰 $d\in[-\bar d,\bar d]$ 想把系统推入。若\textbf{控制更强}（$\bar d<1$）：从危险集外的任何点，控制能以净速度 $1-\bar d>0$ 把系统推开，干扰追不上——有效不安全集不扩张，就是危险集本身 $\{|x|\le1\}$。若\textbf{干扰更强}（$\bar d>1$）：干扰能压制控制，以净速度 $\bar d-1>0$ 把系统拽向危险集——有效不安全集扩张为 $\{|x|\le1+(\bar d-1)|t|\}$。这个对比把“安全集大小取决于控制与干扰的相对操控权”量化了，也预演了 \cref{sec:dg-chase} 追逃博弈“谁操控权强谁占优”的核心。
\end{example}

\subsection{后向可达管（BRT）：在区间内曾经到达}
\label{sec:dg-brt}
安全分析里我们关心的往往不是“恰好在终点出事”，而是“在整个时间窗内\emph{有没有任何一刻}出事”。这就是后向可达\textbf{管}：
\begin{equation}
  \mathcal{R}_{\text{tube}}(t)=\Big\{x:\ \exists\,\gamma,\ \forall a(\cdot),\ \exists\,s\in[t,0],\ \zeta\big(s;x,t,a,\gamma[a]\big)\in\mathcal{T}\Big\}.
  \label{eq:dg-brt-def}
\end{equation}
注意定义里多出来的 $\exists\,s\in[t,0]$——只要\emph{某个时刻}进了目标就算。BRT 对应的方程在 \cref{eq:dg-brs-pde} 上加一个\textbf{冻结项}，保证“一旦进入目标，就永久标记为已到达”，从而让零下水平集随时间单调扩张成一根管。两种等价写法：
\begin{align}
  \min\Big\{D_t V+H(x,\nabla V),\ \ g(x)-V(x,t)\Big\}&=0, \label{eq:dg-brt-vi}\\
  D_t V+\min\big[\,0,\ H(x,\nabla V)\,\big]&=0. \label{eq:dg-brt-mitchell}
\end{align}
\cref{eq:dg-brt-vi} 是变分不等式形式，\cref{eq:dg-brt-mitchell} 是 Mitchell--Tomlin（2005）的原始形式（把哈密顿量整体截断）\cite{mitchell2005time}。两者都实现同一件事：阻止值函数在已进入目标的区域回升，使 $\{V\le0\}$ 只增不减——这就是“管”而非“集”。\cref{ex:dg-integrator} 用的正是 \cref{eq:dg-brt-mitchell}。

\begin{example}[集还是管：让两者分道扬镳的例子]
\label{ex:dg-set-vs-tube}
\cref{ex:dg-integrator} 里 BRS 和 BRT 恰好重合（系统能“走到目标就停下来等”），看不出差别。换一个会\emph{穿过}目标又出来的系统，两者立刻分道扬镳。设想一个钟摆/振子式的状态，在相空间里绕圈：从某个初始状态出发，它的轨迹会在 $t=1\,\mathrm s$ 时\emph{正好扫过}危险集、然后在 $t=2\,\mathrm s$ 时\emph{已经荡出去}、离危险集很远。这个初始状态危险吗？用 \textbf{BRS}（“恰好在终点 $t=2\,\mathrm s$ 是否在危险集内”）判：终点时刻它已荡出去，不在危险集内——BRS \emph{判它安全}。用 \textbf{BRT}（“在 $[0,2\,\mathrm s]$ 内是否曾进入危险集”）判：它在 $t=1\,\mathrm s$ 穿过了危险集——BRT \emph{判它危险}。对安全验证该用 BRT，道理很直接：撞车是“撞过就算撞”，不会因为之后又弹开就当没撞。那个 $\min[0,\cdot]$ 冻结项做的正是这件事——它在轨迹首次进入目标的那一刻把值“锁住”，不让它随后荡出去时再回升到 0 以上。
\end{example}

\begin{insight}
BRS 与 BRT 的全部差别，浓缩在那个冻结项里。BRS（无冻结项，\cref{eq:dg-brs-pde}）回答“恰好在终点能否到达”；BRT（带冻结项，\cref{eq:dg-brt-vi,eq:dg-brt-mitchell}）回答“在窗内曾经能否到达”。安全验证几乎总是要 BRT——因为撞车是“任意时刻撞了都算撞”，不是“只看终点撞没撞”。把 BRS 的方程用到本该用 BRT 的安全分析上，会漏掉那些“中途撞了又弹出去”的危险状态，给出偏乐观、不安全的结论。这就是文献里 $\min[0,\cdot]$ 这个看似神秘的项真正的来历。
\end{insight}

\subsection{reach--avoid：既要到达又要全程避险}
\label{sec:dg-reachavoid}
很多真实任务不是单纯的“到达”或“避免”，而是二者叠加：\textbf{到达一个目标集 $\mathcal{T}$，同时全程不进入危险集 $\mathcal{A}$}。比如无人机要飞到终点、但途中不能撞障碍；自动驾驶要并入目标车道、但全程不能与他车相撞。下面一步步把它的方程搭出来——这能让你看清那个看起来吓人的“双障碍变分不等式”其实只是两件已经会的事拼在一起。

\textbf{第一步：我们已经有了“到达”的零件。}\;前面的 BRT 变分不等式 \cref{eq:dg-brt-vi} 处理纯“到达目标 $\mathcal{T}=\{g\le0\}$”：那个 $\min\{\cdot,\ g-V\}$ 把 $V$ \emph{封顶}在目标距离 $g$ 之下（一旦能到达，$V$ 就掉到 $\le g\le0$，并被冻住不再回升），零下水平集 $\{V\le0\}$ 因此只增不减。记住这个结构：\textbf{内层的 $\min$ + 目标项 $g-V$ = “reach 冻结”}。

\textbf{第二步：加一个“避险”的零件。}\;现在多一个危险集 $\mathcal{A}$，用它的带符号距离 $c(x)$ 编码（$c\le0$ 在 $\mathcal{A}$ 内、$c>0$ 在安全区）。“避险”的要求是：\emph{任何最优路径若必须穿过 $\mathcal{A}$ 才能到达目标，这个起点就不算赢}——哪怕它能到达目标，也要被一票否决。怎么否决？让值函数在危险集内被\emph{强行顶高}到正值（这样它就落不进 $\{V\le0\}$ 这个胜利集）。实现这件事的是一个\textbf{外层的 $\max$ + 障碍项 $-c-V$}：在 $\mathcal{A}$ 内 $c\le0$ 故 $-c\ge0$，外层 $\max\{\cdot,\ -c-V\}=0$ 会强制 $V\ge-c\ge0$，把危险集内的点顶成正值、踢出胜利集；在安全区 $c>0$ 故 $-c<0$，这一项是负的、不起作用，交回内层的 reach 冻结管。记住：\textbf{外层的 $\max$ + 障碍项 $-c-V$ = “avoid 否决”}。

\begin{theorem}[reach--avoid 双障碍变分不等式]
\label{thm:dg-reachavoid}
reach--avoid 的值函数满足“reach 冻结”嵌进“avoid 否决”里的双障碍变分不等式（Margellos--Lygeros 2011\cite{margellos2011hamilton}；Fisac 等 2015 推广到时变情形\cite{fisac2015reach}）：
\begin{equation}
  \max\Big\{\ \underbrace{\min\big\{\,D_t V+H,\ \ g(x)-V\,\big\}}_{\text{reach 冻结（朝目标推进）}},\ \ \underbrace{-c(x)-V}_{\text{avoid 否决（穿障碍即出局）}}\ \Big\}=0,
  \label{eq:dg-reachavoid}
\end{equation}
其中 $g$ 是目标 $\mathcal{T}$ 的带符号距离、$c$ 是危险集 $\mathcal{A}$ 的带符号距离（均用“集合内 $\le0$”约定）；reach--avoid 集（捕获域）就是 $\{x:V(x,t)\le0\}$。$H$ 的 $\max/\min$ 仍按口诀——这里控制争取“存在一条到达且安全的路径”故对控制取 $\min$、干扰处处作对取 $\max$。
\end{theorem}

\textbf{一个一致性检查。}\;若根本没有危险集（$\mathcal{A}=\varnothing$，相当于 $-c\equiv-\infty$），外层那一项永远是负无穷、不起作用，\cref{eq:dg-reachavoid} 立刻塌回纯 reach 冻结 \cref{eq:dg-brt-vi}——这正说明 \cref{eq:dg-reachavoid} 是 BRT 的真正推广，多出来的外层 $\max$ 就是“避险”这一件事。

\begin{insight}
\textbf{双障碍 = 两个方向的“夹”.}\;\cref{eq:dg-reachavoid} 看着复杂，本质是从两个方向夹住值函数。reach 冻结从\emph{上方}把 $V$ 压向目标（$V\le g$，鼓励集合朝目标长大）；avoid 否决从\emph{下方}把危险区的 $V$ 顶起来（$V\ge-c$，把碰障碍的点逐出）。两个障碍一上一下，夹出的零下水平集恰好是“能到达且全程安全”的那批起点。这也是为什么它叫“双障碍（double-obstacle）”——不是两个物理障碍，而是值函数被两个方向的约束同时挡住。Fisac 等用这一个方程，连时变的目标/障碍/动力学都一并处理了，且数值代价与时不变情形几乎一样。
\end{insight}

把本节的三种可达性放到一起，能看清 reach--avoid 在谱系里的位置：\textbf{纯到达}（BRT，“窗内曾到达目标”，冻结项锁成功）、\textbf{纯避险}（把危险集当目标算 BRT、取补集得安全集，“全程未进危险集”）、\textbf{reach--avoid}（既到达又全程避险，\cref{eq:dg-reachavoid} 把到达冻结与避险障碍合进一个方程）。三者共用同一台机器——HJI 演化 + 水平集——区别只在装了几个、哪种约束项。这也是为什么掌握了 BRS/BRT，reach--avoid 只是“再加一个外层 $\max$”而非全新的东西。这是后续处理“带安全约束的交互”时反复要用的形式，也与规划中的安全走廊一脉相承。

\subsection{前向可达、信息优势与通往 CBF}
\label{sec:dg-frs-cbf}
\paragraph{前向可达：从初始集出发能到哪。}
到目前讲的都是“后向”——从目标倒推初始状态。对偶地也有\textbf{前向可达集/管（FRS/FRT）}：给定一个初始状态集 $\mathcal{X}_0$，系统在时长 $|t|$ 内能到达的所有状态（管则是“曾经到达过的”）。前向问题在求解上与后向几乎一样，只是把终端值 PDE 换成初始值 PDE（两者可经时间变量代换互相转化）。用途不同：后向集回答“哪些初始状态危险”（适合安全验证），前向集回答“系统会扩散到哪里”（适合包络分析、碰撞预测、给下游规划划定可行域）。记住一句话区分：\textbf{后向问“从哪来会出事”，前向问“从这出发会到哪”。}

\paragraph{非预期策略：为什么把信息优势给干扰。}
前面 BRS/BRT 的定义里都出现了“$\exists\gamma$”——干扰用一个非预期策略 $\gamma$（\cref{sec:dg-isaacs} 已引入）。为什么安全验证要把这个瞬时信息优势给干扰、而不是给控制？因为安全保证要覆盖\emph{最坏情况}。让干扰“见招拆招”，等于假设了一个最难对付的对手；在这种最悲观的假设下仍判定为安全的状态，才是真的、可信的安全。这与 \cref{prop:dg-maxmin} 的 $\max\min\le\min\max$ 一脉相承：给干扰信息优势取到的是下值，它给出\emph{保守}的安全集。若反过来把信息优势给控制，会高估安全集——把一些其实危险的状态误判为安全。

\paragraph{安全集、最优安全控制与“最小限制滤波器”——通往 CBF。}
可达性分析的产物不止一个集合，还附带一个控制器，这对工程极其有用。当目标 $\mathcal{T}$ 取“危险/不安全状态”、玩家 2 取最坏情况干扰时：BRT $\mathcal{R}_{\text{tube}}(t)$ 是\textbf{有效不安全集}（从中出发，干扰能在窗内某刻把系统逼入危险，无论控制怎么挡）；它的补集是\textbf{安全集}；\cref{eq:dg-brs-ham} 里 $\arg\max_a$ 给出把系统留在安全集内的\textbf{最优安全控制} $a^*(x,t)$。这个安全控制有一个漂亮的“\textbf{最小限制（least-restrictive）}”性质：在安全集\emph{内部}，系统离边界还远，怎么动都安全、控制完全自由；只有当系统逼近安全集\emph{边界}时，$a^*$ 才介入、把系统推回。

\begin{note}
\textbf{多视角理解（HJI 安全滤波 $\leftrightarrow$ CBF 安全滤波）.}\;这个“平时不管、临界才介入”的最小限制思想，正是控制屏障函数（Control Barrier Function, CBF）安全滤波器的灵魂。相似之处：两者都用一个标量函数刻画安全（HJI 用值函数 $V$ 的零下水平集 $\{V\le0\}$ 表示不安全、CBF 用某函数 $h$ 的零上水平集 $\{h\ge0\}$ 表示安全），都用这个函数的梯度构造一个“只在边界处约束控制”的滤波器。事实上，HJI 的安全值函数本身就可以直接当作一个 CBF 用——它的零水平集就是精确的安全边界，$\nabla V$ 给出约束方向。不同之处：HJI 值函数是\emph{算出来的精确}安全证书（但受维度诅咒所限，\cref{sec:dg-curse}），CBF 通常是\emph{手工设计的、可能保守}的安全证书（但便宜、可在线求解）。一句话：CBF 是“廉价、可在线、可能保守”的安全证书，HJI 值函数是“精确但算不动”的安全证书，二者是同一思想在计算代价两端的化身。\pz{控制部深讲：CBF/CLF 的构造、QP 安全滤波、可前向不变性的严格证明留待控制部专章，此处仅给规划/安全验证用途与直觉，不重复长推导。}
\end{note}

\begin{note}
\textbf{多视角理解（BRS/BRT $\leftrightarrow$ Tube MPC 的 RPI 集）.}\;这条联系把博弈和不确定性规划接了起来。不确定性规划里的 Tube MPC，把有界扰动 $w\in\mathcal{W}$ 当作一个对抗玩家，计算一个鲁棒正不变（Robust Positively Invariant, RPI）集——名义轨迹周围的一根“管子”，保证不管扰动怎么取，真实状态都待在管内。这与 HJI 可达性的 BRT 是同一种思想的两个化身：都把最坏情况干扰建模成对手，都算一根“在干扰下系统必然待在/必然落入的区域”的管。相似之处：对手即干扰、产物即不变/可达管、目的都是给出可证明的安全保证。不同之处：Tube MPC 通常在线性/局部、用集合递推（Minkowski 和）高效算，得到不变管；HJI 用 PDE 在网格上算，能处理一般非线性、得到可达管，但代价是维度诅咒（\cref{sec:dg-curse}）。一句话：RPI 集是 BRT 在线性鲁棒控制里的高效特例。\pz{控制部深讲：Tube MPC、RPI 集的 Minkowski 递推与鲁棒 MPC 稳定性证明留待控制部 MPC 章，此处仅建立与 BRT 的概念同源关系。}
\end{note}

\begin{pitfall}[概念误区：混淆后向可达集（BRS）与后向可达管（BRT）]
把“在终点到达”的 BRS 方程（无冻结项）用到本该问“窗内曾经到达”的安全分析上；或反过来。许多人甚至以为带 $\min[0,\cdot]$ 的式子算的是 BRS。\textbf{后果}：用 BRS 做安全验证，会漏掉“中途撞了、终点又恰好弹出危险区”的状态，安全集被高估，得到偏乐观甚至致命的结论。\textbf{根因}：BRS 只看终端时刻是否在目标内；BRT 看整个时间窗内是否曾进入目标。撞车这类安全事件是“任意时刻发生都算”，对应 BRT；$\min[0,\cdot]$ / 变分不等式里的冻结项正是把“集”变成“管”的机制。\textbf{正解}：安全验证一律用 BRT（带冻结项的 \cref{eq:dg-brt-vi,eq:dg-brt-mitchell}）。检验法：问“我关心的坏事是‘只在终点发生才算’，还是‘窗内任何时刻发生都算’？”——后者就必须用 BRT。
\end{pitfall}

\begin{pitfall}[思维陷阱：哈密顿量里 max/min 取反（reach 与 avoid 搞混）]
不区分玩家角色，凭习惯把哈密顿量写成 $\min_a\max_b$ 或 $\max_a\min_b$，赌对一半。\textbf{后果}：算出的可达集张错了方向——本该收缩的在扩张，安全集和危险集对调，可视化出来“看着像那么回事”却完全错误，且极难自查。\textbf{根因}：max/min 的取法由“对每个玩家问存在还是任意”唯一决定，弄反等于把控制和干扰的角色互换，求的是另一个完全不同的博弈。\textbf{正解}：套用“存在取 $\min$、任意取 $\max$”的口诀，先写清可达集定义里每个玩家是 $\exists$ 还是 $\forall$，再据此定哈密顿量。检验法：把干扰强度调到 0，看可达集是否退化为单人最优控制的可达集（如 \cref{ex:dg-integrator} 退化为 $\{|x|\le1+|t|\}$）——若方向不对，多半是 max/min 取反了。
\end{pitfall}

\begin{practice}
\begin{enumerate}
  \item （精读 + 推导）精读 Mitchell--Bayen--Tomlin（2005）的可达性表述部分，自己重走一遍“可达集 $\leftrightarrow$ HJI 零下水平集”的等价论证。重点回答：为什么 BRT 要用 $\min[0,\cdot]$（或等价的变分不等式）这个截断/冻结项，而 BRS 不用？把这与“避免进入”的语义对应起来，并与 \cref{ex:dg-integrator} 的验证对照。
  \item （跨专题思考题）HJI 可达性（零和博弈）和不确定性规划里的 Tube MPC 在概念上是什么关系？请论证“Tube MPC 把有界扰动 $w\in\mathcal{W}$ 当作一个对抗玩家，它的 RPI 不变集 $\approx$ HJI 的后向可达管”，指出这个类比成立的地方，以及它在哪里会失效（提示：从动力学是否线性、集合如何表示两个角度想）。
  \item （角色辨析 + reach--avoid）对“无人机要在窗内\emph{始终}待在一个安全走廊内、同时\emph{最终}到达终点”这一 reach--avoid 任务，写出你会对控制和干扰分别问“存在”还是“任意”，据此定出哈密顿量里每个 $\min/\max$ 的取法。你需要 BRS、BRT，还是 reach--avoid 的组合？说明理由，并解释障碍项在这里起什么作用。
\end{enumerate}
\end{practice}

% ════════════════════════════════════════════════════════════════
\section{水平集数值方法：WENO 与 TVD-RK}
\label{sec:dg-numerics}

\paragraph{动机：从连续 PDE 到网格上的数值演化。}
HJI 方程 \cref{eq:dg-brs-pde} 是连续状态空间上的偏微分方程，计算机没法直接处理连续对象。标准做法是把状态空间划成网格，在每个网格点上存一个值 $V$，把 PDE 的时间演化离散成一步步的更新——从终端条件 $V(x,0)=g(x)$ 出发，沿时间负方向一步步推回去，每一步用相邻网格点的值近似空间梯度 $\nabla V$、算出哈密顿量、更新这一点的值。听起来很直接，难点全在“怎么近似 $\nabla V$”和“怎么走时间步”——因为 \cref{sec:dg-viscosity} 告诉我们值函数有不可微的棱，天真的近似会在棱处崩掉。

\paragraph{反面：天真的中心差分会怎样。}
最自然的想法是用中心差分近似梯度：$\partial_x V\approx(V_{i+1}-V_{i-1})/(2\Delta x)$。在光滑区域这没问题，但在值函数的棱附近，中心差分会产生剧烈的\textbf{伪振荡}：它“看不见”间断的方向性，把棱两侧的信息对称地搅在一起，导致数值解在棱附近上下抖动、长出毛刺，可达集边界变得锯齿化甚至断裂。更糟的是，不带耗散的格式可能收敛到一个\emph{几乎处处满足方程、却不是黏性解}的假解——正是 \cref{sec:dg-viscosity} 里 eikonal 的 $|x|-1$ 那类被否决的解。方程本身选不出正确解，数值格式也一样：格式必须“尊重黏性”，才能收敛到对的那个。

\subsection{尊重黏性的数值格式}
\label{sec:dg-scheme}
\paragraph{数值哈密顿量与迎风思想。}
要让格式稳定且收敛到黏性解，关键是用\textbf{带方向性、带数值耗散}的格式构造一个\textbf{数值哈密顿量} $\hat H$ 去近似 $H(x,\nabla V)$。最基础的是 Lax--Friedrichs 型数值哈密顿量：它在中心差分的基础上人为加一个正比于梯度跳变的耗散项，
\begin{equation}
  \hat H=H\!\left(x,\frac{p^++p^-}{2}\right)-\frac{\alpha}{2}\,(p^+-p^-),
  \label{eq:dg-lf}
\end{equation}
其中 $p^\pm$ 是从两侧估计的梯度、$\alpha$ 是与哈密顿量对梯度的敏感度（最大特征速度）相关的耗散系数。那个 $-\tfrac{\alpha}{2}(p^+-p^-)$ 项就是数值耗散：它在梯度有跳变（棱）的地方自动加阻尼，把伪振荡压下去。

\begin{insight}
\textbf{数值耗散 = 消失黏性的离散身影.}\;这里加的数值耗散项，正是 \cref{sec:dg-viscosity} 里“消失黏性”$\varepsilon\Delta V$ 的离散版本。\cref{deriv:dg-vanishing} 用“加无穷小扩散再取极限”在理论上挑出黏性解；数值格式用“加一点正比于网格的耗散”在计算上逼近同一个黏性解。网格越细，等效的 $\varepsilon$ 越小，数值解越接近真正的黏性解。不同之处：理论里 $\varepsilon\to0$ 是取极限，数值里耗散量由网格步长决定、不能真取 0（否则振荡回来）。这条对应是理解“为什么单调格式收敛到黏性解”的钥匙——收敛到黏性解，不是格式的幸运巧合，而是它内建的数值耗散在兑现 \cref{sec:dg-viscosity} 的黏性极限。
\end{insight}

\paragraph{空间离散：WENO5。}
Lax--Friedrichs 只有一阶精度，在光滑区域太耗散（把本该锐利的特征也磨钝了）。工程上用高阶的\textbf{加权本质无振荡格式（Weighted Essentially Non-Oscillatory, WENO）}——常用工具默认用五阶 WENO（WENO5）。WENO 的核心思想：估计某点梯度时，不固定用一组网格点，而是从几组候选模板里\emph{按光滑度自适应加权}——在光滑区给所有模板高权重（达到高阶精度），在棱附近自动给“跨过间断”的模板几乎零权重（避免把间断信息搅进来）。这就是“本质无振荡”：高阶精度与棱处不振荡两者兼得。把 WENO 的思路拆慢：要点是一个\textbf{两难}——估计梯度时用到的网格点越多（模板越宽），光滑区精度越高（好事）；但模板一旦“跨过”一道棱，就会把棱两侧本不该混的信息搅在一起、产生伪振荡（坏事）。固定用宽模板 $\to$ 高精度但棱处振荡；固定用窄模板 $\to$ 棱处稳但精度低。WENO 怎么两全：

\begin{enumerate}
  \item \textbf{先给出几组候选模板}。对要估梯度的那个点，准备几组不同的相邻网格点组合（偏左、居中、偏右）。每组都能给一个梯度估计，但它们“看”的是不同侧的数据。
  \item \textbf{ENO 的想法——只挑最光滑的一组}。早期 ENO（Essentially Non-Oscillatory）格式：哪组候选模板覆盖的数据最光滑（没跨过棱）就用哪组、丢掉其余。代价：硬性二选一，浪费了其他模板的信息，精度没拉满。
  \item \textbf{WENO 的改进——按光滑度加权混合}。WENO 不做非黑即白的选择，而是给每组模板算一个\emph{光滑度指标}，据此分配权重：光滑的模板给大权重、跨棱的给近乎零的权重，再把各组估计加权平均。于是光滑区所有模板都被“信任”、加权后达到高阶精度（五阶就是 WENO5）；棱附近跨棱的那组被自动压成零权重、退化为只用光滑侧的低阶估计但稳定不振荡。
\end{enumerate}

\begin{note}
\textbf{多视角理解（WENO 的自适应加权 $\leftrightarrow$ 鲁棒统计里的加权）.}\;WENO“按光滑度给模板分权重、把异常的那组压到近零”这招，和鲁棒统计里“按残差给样本分权重、把离群点压低”是同一种智慧。相似之处：都在“我有几个来源、其中某些被污染（跨棱/离群），该信谁多一点”的问题上，用一个连续权重而非硬性取舍来兼顾精度与稳健。不同之处：WENO 的“污染”是几何的（模板跨过间断）、权重由光滑度指标定；鲁棒统计的“污染”是数据的、权重由残差定。看懂这个类比，就不会把 WENO 当成一堆神秘系数，而是“对若干梯度估计做稳健加权融合”。
\end{note}

\paragraph{时间推进：TVD-RK3。}
空间离散后，HJI 变成一个关于网格值向量的常微分方程组，需要时间积分器往回推。直接用高阶 Runge--Kutta 可能引入振荡，所以用\textbf{总变差不增（Total Variation Diminishing, TVD）的三阶 Runge--Kutta（TVD-RK3，又称 SSP-RK3）}。它的设计保证：只要每个子步是稳定的（不增加解的总变差），三阶组合后整体也不增加总变差——即时间方向上也“不无中生有地长出振荡”，与空间的 WENO 配合，全程保持非振荡。

\paragraph{CFL 条件：时间步不能太大。}
显式时间推进有个硬约束——\textbf{CFL（Courant--Friedrichs--Lewy）条件}：一个时间步内，信息在数值上传播的距离不能超过一个网格格子，否则格式发散。形式上 $\Delta t\lesssim\Delta x/\alpha$（$\alpha$ 是最大特征速度）。直觉：物理信息以有限速度传播，数值格式每步只能“看”到相邻格点，若 $\Delta t$ 太大、真实信息在一步内越过了好几个格子，格式就跟丢了、爆掉。所以网格越细（$\Delta x$ 越小），允许的 $\Delta t$ 越小、步数越多——这是细网格代价的一部分。

\subsection{从零写一个最小可达性求解器}
\label{sec:dg-mincode}
成熟工具好用，但“会调库”不等于“懂数值”。下面用纯 NumPy、零依赖从头写一个一维可达性求解器，把 Lax--Friedrichs 数值哈密顿量、CFL、冻结项全落成几行能跑的代码，并用 \cref{ex:dg-integrator} 的解析解 BRT $=\{|x|\le1+T\}$ 检验它对不对。对象就是那个单积分器 $\dot x=u,\ u\in[-1,1]$，目标 $\{|x|\le1\}$，控制争取到达。BRT 的“时间-到-达”形式是 $V_\tau=-|V_x|$（零下水平集以单位速度向外扩）。

\begin{codebox}[Python]{最小一维可达性求解器（Lax--Friedrichs + CFL）}
import numpy as np

# 1D single integrator x_dot = u, u in [-1,1]; target {|x|<=1}; compute T-horizon BRT.
# Time-to-reach form of HJI(BRT): V_tau = -|V_x|  ({V<=0} expands at unit speed).
xs = np.linspace(-6.0, 6.0, 1201)
dx = xs[1] - xs[0]
V  = np.abs(xs) - 1.0                  # init = signed distance g(x)=|x|-1 (Sec reach)
T, alpha = 3.0, 1.0                    # horizon; alpha = max|dH/dp| = 1 (max char speed)
dt = 0.4 * dx / alpha                  # CFL: dt <~ dx/alpha, 0.4 safety margin
for _ in range(int(T / dt)):
    pm = np.empty_like(V); pp = np.empty_like(V)
    pm[1:]  = (V[1:] - V[:-1]) / dx; pm[0]  = pm[1]     # left diff p-
    pp[:-1] = (V[1:] - V[:-1]) / dx; pp[-1] = pp[-2]    # right diff p+
    pc = 0.5 * (pm + pp)                                # central gradient
    # Lax-Friedrichs numerical Hamiltonian: |V_x|_LF = |pc| - (alpha/2)(pp - pm)
    #   the last term is numerical dissipation = discrete image of vanishing viscosity
    Vx_LF = np.abs(pc) - 0.5 * alpha * (pp - pm)
    V = V - dt * Vx_LF                                  # V_tau = -|V_x|, advance one step
    # note: here H = -|V_x| <= 0, so the (1.4b) min[0,H] freeze term never activates;
    #       in an avoid (safety) problem the freeze term keeps the safe set from shrinking.
lo, hi = xs[V <= 0].min(), xs[V <= 0].max()
print(f"numeric BRT ~ [{lo:.2f}, {hi:.2f}]; analytic = [-{1+T:.0f}, {1+T:.0f}]")
# typical output: numeric BRT ~ [-3.99, 3.99]; analytic = [-4, 4]
\end{codebox}

短短二十几行，把数值三件套都跑了一遍：中心梯度 + 那个 $-\tfrac{\alpha}{2}(p^+-p^-)$ 耗散项就是 Lax--Friedrichs 数值哈密顿量（在 $x=0$ 的尖点处加阻尼，对应 \cref{sec:dg-viscosity} 的消失黏性）；\verb@dt = 0.4*dx/alpha@ 就是 CFL；注释里点明的 \verb@min[0,H]@ 冻结项在这个 reach 例子里恰好不激活，但换成 avoid（安全）问题它就是阻止安全集回缩的关键。跑出来的数值 BRT $[-3.99,3.99]$ 与解析解 $[-4,4]$ 吻合到网格精度——这既验证了代码，也验证了 \cref{ex:dg-integrator} 那个“$\{|x|\le1+T\}$”的手算。

\subsection{用成熟工具实跑一次：Air3D 的 BRT}
\label{sec:dg-tool}
手写版把原理讲透，成熟工具则处理高维、高阶、GPU 加速。本领域最现代、最干净的开源工具之一是 hj\_reachability（JAX 实现），专门求解零和微分博弈的 HJI-PDE，把可达集表示成值函数黏性解的零下水平集——和本章的表述逐字对应。下面用它算 \cref{sec:dg-chase} 那个 Air3D 追逃的后向可达管，顺带把 \cref{sec:dg-scheme} 的概念落到 API 上（接口以工具当前官方文档为准）。

\begin{codebox}[Python]{用 hj\_reachability 求 Air3D 的 BRT（API 概念示意）}
import jax.numpy as jnp
import hj_reachability as hj

# 1) Dynamics: Air3D (built-in). Pursuer/evader fly at fixed speed, bounded turn rate.
#    Relative coords (x, y, psi): the 3D relative dynamics, avoids the absolute 6D state.
dynamics = hj.systems.Air3d(
    evader_speed=5.0, pursuer_speed=5.0,
    evader_max_turn_rate=1.0, pursuer_max_turn_rate=1.0)

# 2) Grid: uniform lattice over relative state. Dim 2 (psi) is angular -> periodic.
grid = hj.Grid.from_lattice_parameters_and_boundary_conditions(
    hj.sets.Box(lo=jnp.array([-6., -10., 0.]),
                hi=jnp.array([20.,  10., 2 * jnp.pi])),
    (51, 41, 50), periodic_dims=2)

# 3) init = signed distance to collision disk {x^2+y^2<=R^2} (level-set encoding)
R = 5.0
initial_values = jnp.linalg.norm(grid.states[..., :2], axis=-1) - R

# 4) "very_high" accuracy ~ high-order WENO spatial discretization;
#    backwards_reachable_tube post-processor = the min[0,.] freeze term -> set becomes tube.
solver_settings = hj.SolverSettings.with_accuracy(
    "very_high",
    hamiltonian_postprocessor=hj.solver.backwards_reachable_tube)

# 5) integrate backward from t=0 to t=-T (TVD-RK inside, auto-satisfies CFL)
T = 2.8
target_values = hj.step(solver_settings, dynamics, grid, 0.0, initial_values, -T)
# BRT (relative configs the evader cannot avoid capture from) = {x : target_values(x) <= 0}
\end{codebox}

把这段与原理对一遍，能看清“理论概念 $\leftrightarrow$ 库参数”的一一对应（\cref{tab:dg-api}）。

\begin{table}[H]\centering\small
\caption{hj\_reachability 参数与本章概念的对应}
\label{tab:dg-api}
\begin{tabular}{p{0.46\textwidth} p{0.46\textwidth}}
\toprule
代码里的东西 & 对应本章的概念 \\
\midrule
\texttt{initial\_values = $\lVert\cdot\rVert$ - R} & \cref{sec:dg-reach} 把目标集编码成带符号距离的零下水平集 \\
\texttt{with\_accuracy("very\_high")} & \cref{sec:dg-scheme} 的高阶 WENO 空间离散（非振荡、棱处不冒毛刺） \\
\texttt{backwards\_reachable\_tube} 后处理 & \cref{eq:dg-brt-mitchell} 的 $\min[0,\cdot]$ 冻结项——选 BRT（管）而非 BRS（集） \\
\texttt{periodic\_dims=2} & \cref{sec:dg-chase} 相对朝向 $\psi$ 是角度、绕一圈相接 \\
\texttt{hj.\allowbreak step(..., 0.0, ..., -T)} & 从终端沿时间负方向积分（TVD-RK 内部保证 CFL） \\
\texttt{target\_values $\le$ 0} & 取零下水平集得到可达管 \\
\bottomrule
\end{tabular}
\end{table}

\begin{insight}
\textbf{理论与工具的同构.}\;一个好的可达性库不是把理论“封装掉”，而是把理论的每个部件\emph{暴露成一个可调参数}——你选的精度档对应 WENO 的阶、你挑的后处理对应 BRS/BRT 的冻结项、你给的初值对应目标集的水平函数。看懂这段代码的前提，恰恰是看懂 \cref{sec:dg-reach,sec:dg-scheme}；反过来，跑通这段代码也最能巩固那些抽象概念。这也是为什么本章先讲透原理、再给工具：工具是原理的接口，不是替代品。
\end{insight}

\subsection{工程权衡与维度诅咒的预兆}
\label{sec:dg-tradeoff}
水平集方法精确、能处理一般非线性、自带安全控制器——但它在网格上算，代价直接由网格规模决定（\cref{tab:dg-knob}）。前两行是常规的“精度换算力”权衡。真正的拦路虎是最后一行的\textbf{维数 $d$}：网格点数 $N^d$ 随维数指数增长，$N=50$、$d=6$ 就是 $50^6\approx1.5\times10^{10}$ 个点，内存与时间都到了单机极限。这就是网格法被卡在约 6 维的根本原因，也正是 \cref{sec:dg-curse} 要正面处理的\textbf{维度诅咒}——以及 \cref{sec:dg-lq} 那条“退而用局部 LQ 近似”路线之所以必要的理由。

\begin{table}[H]\centering\small
\caption{水平集求解的工程旋钮}
\label{tab:dg-knob}
\begin{tabular}{p{0.30\textwidth} p{0.30\textwidth} p{0.32\textwidth}}
\toprule
旋钮 & 调高的收益 & 调高的代价 \\
\midrule
空间分辨率（每维格点数 $N$） & 可达集边界更精确、棱更锐利 & 内存与每步计算量按 $N^d$（$d$ 为维数）增长 \\
时间精度（步数） & 演化更准、CFL 更安全 & 步数随 $\Delta t$ 减小而线性增加 \\
维数 $d$ & 能处理更复杂系统 & \textbf{指数爆炸}：$N^d$ 个网格点 \\
\bottomrule
\end{tabular}
\end{table}

\begin{pitfall}[概念误区：用不带耗散的格式（如纯中心差分）解 HJI]
图省事或图“高精度”，用中心差分近似梯度、用普通 RK 推进时间。\textbf{后果}：在值函数的棱附近出现伪振荡，可达集边界长毛刺、断裂；严重时收敛到一个几乎处处满足方程、却不是黏性解的假解，安全判断完全错误。\textbf{根因}：HJI 值函数有不可微的棱（\cref{sec:dg-viscosity}），不带方向性/耗散的格式无法正确处理间断，且不“尊重黏性”。\textbf{正解}：用为黏性解设计的单调/本质无振荡格式——WENO5（空间）+ TVD-RK3（时间）+ 合适的数值哈密顿量（如 Lax--Friedrichs 型耗散）。检验法：加密网格，看可达集边界是否收敛、是否仍光滑无毛刺。
\end{pitfall}

\begin{pitfall}[思维陷阱：违反 CFL 条件，时间步取太大]
为了少跑几步、加快计算，把 $\Delta t$ 设得很大。\textbf{后果}：数值解发散——值函数出现爆炸性的大数甚至 NaN，可达集面目全非。\textbf{根因}：显式格式每步只能传播一个格子的信息，$\Delta t$ 太大时真实信息越过多个格子，格式跟丢、不稳定。\textbf{正解}：让 $\Delta t\lesssim\Delta x/\alpha$（$\alpha$ 为最大特征速度），网格加密时同步减小 $\Delta t$；多数工具会自动按 CFL 选步长，自己写格式时务必加这个约束。检验法：若细化网格后突然发散，先查 $\Delta t$ 是否随 $\Delta x$ 同步减小。
\end{pitfall}

\begin{practice}
\begin{enumerate}
  \item （概念联系，回扣 \cref{sec:dg-viscosity}）用自己的话说明：为什么“加数值耗散的单调格式”会收敛到黏性解，而中心差分不会？把它和“消失黏性 $\varepsilon\Delta V$、取 $\varepsilon\to0$”对应起来——数值里的 $\varepsilon$ 大致相当于什么量？
  \item （估算维度诅咒）估算用每维 50 格点的均匀网格存储一个 6 维 HJI 值函数需要多少个网格点、多少内存（设每点一个 8 字节 double）。若提到 8 维呢？由此说明为什么网格法卡在约 6 维。
  \item （数值实验）用任一可达性工具对一个 2D 系统（如 \cref{ex:dg-integrator} 的单积分器升到 2D，或 \cref{sec:dg-chase} 的 Air3D）算 BRT。把空间分辨率从粗到细跑几档，观察可达集边界如何收敛；再故意把 $\Delta t$ 调大违反 CFL，记录发散现象。
\end{enumerate}
\end{practice}

% ════════════════════════════════════════════════════════════════
\section{LQ 博弈与耦合 Riccati 方程组}
\label{sec:dg-lq}

\paragraph{动机：唯一能闭式求解的博弈，也是局部近似的归宿。}
一般非线性微分博弈没有解析解，只能靠 \cref{sec:dg-numerics} 的网格法硬算，又被维度诅咒卡死。但有一个特例能闭式求解——\textbf{线性二次（LQ）博弈}：动力学线性、每人代价二次。它之所以重要，有两层：第一，它本身可解，给出反馈 Nash 均衡的闭式表达。第二，也是更深的一层——\textbf{任何光滑的微分博弈，都能在一条名义轨迹附近局部近似成一个 LQ 博弈}（动力学一阶 Taylor 展开成线性、代价二阶展开成二次）。这正是实时博弈求解器 iLQGames 的核心套路：反复地“线性化—二次化—解一个 LQ 博弈—更新轨迹”，直到收敛。所以本节解的这组耦合 Riccati，就是 iLQGames 每一步迭代里那个被反复求解的内核。打通了它，后续的实时博弈求解才有根。提醒一句：与本章前几节的零和设定不同，LQ 博弈这里是\textbf{一般和}——每个玩家有自己的 $\mathbf Q_i,\mathbf R_i$，目标不必对立。这正是 \cref{sec:dg-why} 预告过的“本章中的一般和例外”。

\paragraph{反面：多人 LQ 就是各自解一个 LQR？}
一个很自然、却错得很彻底的想法是：“既然每人都是线性二次，那让每个玩家各自解一个标准 LQR、各拿各的反馈增益不就行了？”这忽略了博弈的耦合：玩家 $i$ 面对的“系统”，其动力学里含着别人的控制 $\mathbf u_j$——而别人的 $\mathbf u_j$ 又是别人反馈增益的产物。你按“别人不动”去解自己的 LQR，得到的增益在别人也按最优反馈行动时\textbf{根本不是最优响应}，于是它不构成纳什均衡。正确的做法必须把所有玩家的最优性条件\emph{联立}求解——这就逼出了“耦合”的 Riccati。

\subsection{N 人 LQ 博弈与反馈 Nash}
\label{sec:dg-lq-setup}
\paragraph{问题设定。}
$N$ 个玩家共享线性动力学，每人有自己的控制输入：
\begin{equation}
  \dot{\mathbf x}=\mathbf A\mathbf x+\sum_{j=1}^N \mathbf B_j \mathbf u_j,\qquad \mathbf x\in\mathbb{R}^n,\ \mathbf u_j\in\mathbb{R}^{m_j}.
  \label{eq:dg-lq-dyn}
\end{equation}
玩家 $i$ 最小化自己的二次代价（这里取只惩罚自身控制的常见形式）：
\begin{equation}
  J_i=\frac12\int_0^T\Big(\mathbf x^\top \mathbf Q_i \mathbf x+\mathbf u_i^\top \mathbf R_{ii}\mathbf u_i\Big)\,\mathrm dt+\frac12 \mathbf x(T)^\top \mathbf Q_i^f \mathbf x(T),\qquad \mathbf Q_i\succeq0,\ \mathbf R_{ii}\succ0.
  \label{eq:dg-lq-cost}
\end{equation}

\paragraph{反馈 Nash 均衡。}
我们找\textbf{反馈}策略 $\mathbf u_i=-\mathbf K_i(t)\,\mathbf x$（依赖当前状态，不是开环的时间函数）——\cref{sec:dg-isaacs} 强调过反馈策略有天然的扰动拒绝能力，也是 iLQGames 输出的形式。“反馈 Nash”的含义：在所有其他玩家的反馈增益 $\mathbf K_j\,(j\ne i)$ 固定时，$\mathbf K_i$ 是玩家 $i$ 的最优反馈。每个玩家的值函数取二次形式 $V_i(\mathbf x,t)=\tfrac12 \mathbf x^\top \mathbf P_i(t)\,\mathbf x$，$\mathbf P_i\succeq0$ 待定。

\begin{derivation}[耦合 Riccati 的推导]
\label{deriv:dg-coupled}
固定其他玩家的策略 $\mathbf u_j=-\mathbf K_j\mathbf x\,(j\ne i)$，玩家 $i$ 面对的是一个\emph{修正过的}线性系统：
\[
  \dot{\mathbf x}=\underbrace{\Big(\mathbf A-\sum_{j\ne i}\mathbf B_j\mathbf K_j\Big)}_{=:\ \tilde{\mathbf A}_i}\mathbf x+\mathbf B_i \mathbf u_i.
\]
这下玩家 $i$ 的问题退化成一个标准单人 LQR——系统矩阵 $\tilde{\mathbf A}_i$、输入矩阵 $\mathbf B_i$、代价权重 $\mathbf Q_i,\mathbf R_{ii}$。单人 LQR 的解（前置桥接 \cref{eq:dg-riccati}）：最优反馈 $\mathbf K_i=\mathbf R_{ii}^{-1}\mathbf B_i^\top \mathbf P_i$，其中 $\mathbf P_i$ 解关于 $(\tilde{\mathbf A}_i,\mathbf B_i,\mathbf Q_i,\mathbf R_{ii})$ 的 Riccati 方程
\[
  -\dot{\mathbf P}_i=\tilde{\mathbf A}_i^\top \mathbf P_i+\mathbf P_i\tilde{\mathbf A}_i+\mathbf Q_i-\mathbf P_i\mathbf B_i\mathbf R_{ii}^{-1}\mathbf B_i^\top \mathbf P_i.
\]
现在把 $\tilde{\mathbf A}_i=\mathbf A-\sum_{j\ne i}\mathbf B_j\mathbf K_j=\mathbf A-\sum_{j\ne i}\mathbf S_j\mathbf P_j$ 代回（记 $\mathbf S_j:=\mathbf B_j\mathbf R_{jj}^{-1}\mathbf B_j^\top$，对称半正定），展开那两个含 $\tilde{\mathbf A}_i$ 的项：
\[
  \tilde{\mathbf A}_i^\top \mathbf P_i+\mathbf P_i\tilde{\mathbf A}_i=\mathbf A^\top \mathbf P_i+\mathbf P_i\mathbf A-\sum_{j\ne i}\big(\mathbf P_j\mathbf S_j\mathbf P_i+\mathbf P_i\mathbf S_j\mathbf P_j\big).
\]
代回整理即得玩家 $i$ 的耦合 Riccati 方程。
\end{derivation}

\begin{theorem}[耦合 Riccati 方程组（反馈 Nash，有限时域）]
\label{thm:dg-coupled}
$N$ 人 LQ 博弈 \cref{eq:dg-lq-dyn,eq:dg-lq-cost} 的反馈 Nash 均衡由下列耦合 Riccati 方程组刻画（$i=1,\dots,N$）：
\begin{align}
  -\dot{\mathbf P}_i&=\mathbf A^\top \mathbf P_i+\mathbf P_i\mathbf A+\mathbf Q_i-\mathbf P_i\mathbf S_i\mathbf P_i-\sum_{j\ne i}\big(\mathbf P_j\mathbf S_j\mathbf P_i+\mathbf P_i\mathbf S_j\mathbf P_j\big),\qquad \mathbf P_i(T)=\mathbf Q_i^f, \label{eq:dg-coupled}\\
  \mathbf S_j&=\mathbf B_j\mathbf R_{jj}^{-1}\mathbf B_j^\top,\qquad \mathbf K_i=\mathbf R_{ii}^{-1}\mathbf B_i^\top \mathbf P_i. \label{eq:dg-gain}
\end{align}
闭环动力学为 $\dot{\mathbf x}=\big(\mathbf A-\sum_j \mathbf S_j\mathbf P_j\big)\mathbf x$。从终端 $\mathbf P_i(T)=\mathbf Q_i^f$ 出发，沿时间负方向\emph{同时}反向积分这 $N$ 条方程，每一步每条方程都要用到所有其他 $\mathbf P_j$。
\end{theorem}

\paragraph{为什么“耦合”。}
看 \cref{eq:dg-coupled} 里那个求和项 $-\sum_{j\ne i}(\mathbf P_j\mathbf S_j\mathbf P_i+\mathbf P_i\mathbf S_j\mathbf P_j)$——它把玩家 $i$ 的 $\mathbf P_i$ 和所有别人的 $\mathbf P_j$ 缠在了一起。玩家 $i$ 的最优反馈依赖别人的反馈（通过 $\tilde{\mathbf A}_i$ 里的 $\sum_{j\ne i}\mathbf S_j\mathbf P_j$），别人的又依赖 $i$ 的，循环往复。所以这 $N$ 条方程不能各解各的，必须当成一个整体同时求解。

\begin{derivation}[把两人情形写全：交叉项逐项现身]
\label{deriv:dg-2p}
一般式 \cref{eq:dg-coupled} 看着密，把它落到 $N=2$ 写全，就能看清交叉项是怎么从单人 Riccati 里“长”出来的：
\begin{align*}
  -\dot{\mathbf P}_1&=\underbrace{\mathbf A^\top \mathbf P_1+\mathbf P_1\mathbf A+\mathbf Q_1-\mathbf P_1\mathbf S_1\mathbf P_1}_{\text{玩家 1 单独存在时的标准 Riccati}}-\underbrace{(\mathbf P_2\mathbf S_2\mathbf P_1+\mathbf P_1\mathbf S_2\mathbf P_2)}_{\text{玩家 2 介入带来的修正}},\\
  -\dot{\mathbf P}_2&=\underbrace{\mathbf A^\top \mathbf P_2+\mathbf P_2\mathbf A+\mathbf Q_2-\mathbf P_2\mathbf S_2\mathbf P_2}_{\text{玩家 2 单独存在时的标准 Riccati}}-\underbrace{(\mathbf P_1\mathbf S_1\mathbf P_2+\mathbf P_2\mathbf S_1\mathbf P_1)}_{\text{玩家 1 介入带来的修正}}.
\end{align*}
把第一条逐段读：前半截 $\mathbf A^\top \mathbf P_1+\mathbf P_1\mathbf A+\mathbf Q_1-\mathbf P_1\mathbf S_1\mathbf P_1$ \emph{就是}前置桥接里那条标准单人 Riccati——如果玩家 2 不存在，玩家 1 解的恰好是它。后半截 $-(\mathbf P_2\mathbf S_2\mathbf P_1+\mathbf P_1\mathbf S_2\mathbf P_2)$ 从哪来？回到推导：玩家 1 面对的不是原系统 $\mathbf A$，而是被玩家 2 反馈改写过的 $\tilde{\mathbf A}_1=\mathbf A-\mathbf S_2\mathbf P_2$（这里 $\mathbf S_2\mathbf P_2=\mathbf B_2\mathbf K_2$ 正是玩家 2 的反馈作用）。把 $\tilde{\mathbf A}_1$ 代进 $\tilde{\mathbf A}_1^\top \mathbf P_1+\mathbf P_1\tilde{\mathbf A}_1$ 并减去原来的 $\mathbf A^\top \mathbf P_1+\mathbf P_1\mathbf A$，多出来的正是
\[
  -(\mathbf S_2\mathbf P_2)^\top \mathbf P_1-\mathbf P_1(\mathbf S_2\mathbf P_2)=-(\mathbf P_2\mathbf S_2\mathbf P_1+\mathbf P_1\mathbf S_2\mathbf P_2),
\]
（用了 $\mathbf S_2$ 对称、$(\mathbf S_2\mathbf P_2)^\top=\mathbf P_2\mathbf S_2$）——这就是那个交叉项的全部来历。一句话：交叉项 = “对手的反馈改写了我的系统矩阵”留下的痕迹。
\end{derivation}

\begin{insight}
耦合 Riccati = $N$ 条单人 Riccati 被那些交叉项 $\mathbf P_j\mathbf S_j\mathbf P_i$ 相互纠缠在一起。把交叉项去掉，\cref{eq:dg-coupled} 就退化成 $N$ 条互不相干的标准 Riccati——那对应“每人当别人不存在、各解各的 LQR”的错误做法。正是那些交叉项，把“各自最优”修正成了“互为最优响应”。换句话说，这个纠缠就是“纳什均衡”在线性世界里的数学化身——均衡的本质（我的最优依赖你的最优）被精确地编码进了交叉项。
\end{insight}

\begin{note}
\textbf{对比性思维（单人 Riccati vs N 人耦合 Riccati）.}\;单人 LQR 的 Riccati $-\dot{\mathbf P}=\mathbf A^\top \mathbf P+\mathbf P\mathbf A+\mathbf Q-\mathbf P\mathbf S\mathbf P$，一条方程、自洽、在可镇定性下必有解。$N$ 人耦合 Riccati \cref{eq:dg-coupled} 多了交叉项，$N$ 条方程相互依赖、必须联立。差别就在这些交叉项上——它们是“博弈”相对“最优控制”的全部线性增量。一个关键后果：单人 Riccati 的解在温和条件下保证存在，\textbf{耦合 Riccati 的解却不保证存在}。交叉项可能导致反向积分在有限时间内发散（finite escape），即便每个玩家单独的 LQR 都好端端有解。所以用 \cref{eq:dg-coupled} 时必须检查解是否存在、是否在整个时域上有界，不能想当然。
\end{note}

\begin{example}[两人标量 LQ 博弈，含解析解]
\label{ex:dg-scalar-lq}
把 \cref{eq:dg-coupled} 落到最小的可解例子。设标量状态 $x\in\mathbb{R}$，两个玩家 $\dot x=ax+b_1u_1+b_2u_2$，$J_i=\frac12\int_0^\infty(q_i x^2+r_i u_i^2)\,\mathrm dt$。记 $s_i=b_i^2/r_i$。看\emph{稳态}（无穷时域，令 $\dot P_i=0$，$P_i\to p_i$ 为标量常数），\cref{eq:dg-coupled} 化为两条耦合代数方程
\[
  0=2a\,p_1+q_1-s_1p_1^2-2s_2p_2p_1,\qquad 0=2a\,p_2+q_2-s_2p_2^2-2s_1p_1p_2.
\]
取\emph{对称情形}（$q_1=q_2=q$，$s_1=s_2=s$，由对称性 $p_1=p_2=p$）一条方程就够：
\[
  0=2a\,p+q-sp^2-2sp^2=2a\,p+q-3sp^2\ \Longrightarrow\ 3sp^2-2ap-q=0,
\]
解得（取正根）$p=\dfrac{a+\sqrt{a^2+3sq}}{3s}$。对照\emph{单人 LQR}（玩家 1 当玩家 2 不存在，去掉交叉项 $2s_2p_2p_1$）：$0=2ap+q-sp^2\Rightarrow p=\dfrac{a+\sqrt{a^2+sq}}{s}$。两者明显不同：博弈解的分母是 $3s$、根号里是 $a^2+3sq$，比单人解小。物理含义：当两个玩家都用各自的控制影响同一个状态、各自承担二次代价时，每人“愿意”施加的控制力度（增益 $\propto p$）比“独占系统”时更克制——因为对方也在使劲，过度用力既浪费自己的代价、又被对方部分抵消。均衡是双方“互相迁就”后的不动点，而非各自的最优。存在性方面：此例判别式 $a^2+3sq>0$ 恒成立（$s,q>0$），故实解总存在；但换成非对称权重或有限时域、特定 $\mathbf Q_i$，耦合方程可能无解或反向积分发散——这正是上面强调的、单人 LQR 没有的麻烦。
\end{example}

\subsection{参考实现：耦合 Riccati 反向积分器}
\label{sec:dg-lq-code}
把 \cref{eq:dg-coupled} 写成能跑的代码，是理解 iLQGames 内核最直接的方式。下面给一个自包含、可直接运行的两人版参考实现（NumPy + SciPy），关键行注明对应公式哪一项。

\begin{codebox}[Python]{两人 LQ 博弈耦合 Riccati 反向积分（公式 dg-coupled / dg-gain）}
import numpy as np
from scipy.integrate import solve_ivp

def coupled_riccati_2p(A, B1, B2, Q1, Q2, R1, R2, Qf1, Qf2, T, N=400):
    """Feedback-Nash coupled Riccati for a 2-player finite-horizon LQ game.
    dynamics: x_dot = A x + B1 u1 + B2 u2
    cost  i:  Ji = int_0^T (x'Qi x + ui'Ri ui) dt + x(T)'Qfi x(T)
    returns:  real-time grid ts, P1(t)/P2(t), gains K1(t)/K2(t) (policy ui = -Ki x)
    """
    n = A.shape[0]
    S1 = B1 @ np.linalg.solve(R1, B1.T)     # S_j = B_j R_jj^{-1} B_j' (sym PSD)
    S2 = B2 @ np.linalg.solve(R2, B2.T)

    def deriv(tau, y):
        # tau = T - t turns backward integration into a forward one: dPi/dtau = -Pi_dot = RHSi
        P1 = y[:n * n].reshape(n, n)
        P2 = y[n * n:].reshape(n, n)
        # RHS1 = A'P1 + P1 A + Q1 - P1 S1 P1 - (P2 S2 P1 + P1 S2 P2)   <- player 1
        dP1 = A.T @ P1 + P1 @ A + Q1 - P1 @ S1 @ P1 - (P2 @ S2 @ P1 + P1 @ S2 @ P2)
        # RHS2 = A'P2 + P2 A + Q2 - P2 S2 P2 - (P1 S1 P2 + P2 S1 P1)   <- player 2
        dP2 = A.T @ P2 + P2 @ A + Q2 - P2 @ S2 @ P2 - (P1 @ S1 @ P2 + P2 @ S1 @ P1)
        return np.concatenate([dP1.ravel(), dP2.ravel()])

    y0 = np.concatenate([Qf1.ravel(), Qf2.ravel()])    # tau=0 <-> real terminal t=T, Pi(T)=Qfi
    taus = np.linspace(0.0, T, N)
    sol = solve_ivp(deriv, (0.0, T), y0, t_eval=taus, rtol=1e-9, atol=1e-11)
    if not sol.success:                                # backward divergence = unbounded solution
        raise RuntimeError("coupled Riccati diverged: solution may not exist (pitfall)")

    ts = T - taus                                      # back to real time (decreasing T -> 0)
    P1 = sol.y[:n * n].T.reshape(-1, n, n)
    P2 = sol.y[n * n:].T.reshape(-1, n, n)
    K1 = np.stack([np.linalg.solve(R1, B1.T @ P) for P in P1])   # Ki = Ri^{-1} Bi' Pi
    K2 = np.stack([np.linalg.solve(R2, B2.T @ P) for P in P2])
    return ts, P1, P2, K1, K2
\end{codebox}

几处值得停下来看的设计：\textbf{反向变正向}——Riccati 是终端值问题（给 $\mathbf P_i(T)$ 倒推），代码用 $\tau=T-t$ 的换元把它变成从 $\tau=0$ 正向积分（初值就是终端权重 $\mathbf Q_i^f$），于是能直接喂给标准正向积分器；\textbf{交叉项就是博弈}——\verb@deriv@ 里 \verb@-(P2 S2 P1 + P1 S2 P2)@ 这一项，正是上面本质洞察里说的“把两人的 $\mathbf P$ 纠缠起来”的耦合项，删掉它两条方程立刻解耦成各自独立的单人 Riccati；\textbf{发散即无解}——\verb@if not sol.success@ 那行直接对应耦合 Riccati 可能在有限时间内发散（finite escape）。\textbf{验证纳什}：拿到 $\mathbf K_1,\mathbf K_2$ 后，固定玩家 2 的增益、对修正系统 $\tilde{\mathbf A}=\mathbf A-\mathbf B_2\mathbf K_2$ 重解玩家 1 的\emph{单人} LQR（标准 CARE），应当得回同一条 $\mathbf K_1$——得回了才确认是纳什均衡而非两个各自为政的 LQR。其背后数学正是 \cref{deriv:dg-coupled} 本身：玩家 1 面对的修正系统 $\tilde{\mathbf A}=\mathbf A-\mathbf S_2\mathbf P_2$ 对它解单人 LQR 的代数 Riccati，展开后与耦合方程 \cref{eq:dg-coupled} 的玩家 1 那条逐项相同。\pz{存疑：源文档报告对称标量例 $a=0.5,b_1=b_2=1,q=1,r=1$（$s=1$）下博弈解 $P_1(0)=0.7676$、单人 LQR $=1.6180$、纳什自检误差约 $4.3\times10^{-10}$；解析式 $\tfrac{a+\sqrt{a^2+3sq}}{3s}$ 与 $\tfrac{a+\sqrt{a^2+sq}}{s}$ 代入确给 0.7676 与 1.6180，但具体数值结果未独立复跑，待联网/运行核实。}

\begin{note}
\textbf{理论--工程桥接（这段代码就是 iLQGames 的内核）.}\;这个 \verb@coupled_riccati_2p@ 函数，正是 iLQGames 每一步迭代反复调用的那块。iLQGames 处理一般非线性博弈，但它每一步做的事就是：沿名义轨迹把动力学线性化成 $\mathbf A_k,\mathbf B_{i,k}$、代价二次化成 $\mathbf Q_{i,k}$，然后\emph{解一次本函数所做的耦合 Riccati}，拿到反馈增益去更新轨迹。也就是说——iLQGames = 在每个时刻把一般博弈近似成本节的 LQ 博弈，反复解耦合 Riccati 直到收敛。读懂这段，就读懂了实时博弈求解器跳动的心脏。顺带一提，耦合 Riccati 的反向递推在离散时间下呈块三对角结构，与 SLAM 里因子图的稀疏 Hessian 同构。\pz{存疑：源文档称“iLQGames.jl 约 70 行核心代码”“C++ ilqgames 的 lq\_solver.cc 约 200 行 Eigen”，行数为约数、未核，仅作量级直觉。}
\end{note}

\subsection{反馈 Nash vs 开环 Nash}
\label{sec:dg-nash-types}
\cref{eq:dg-coupled} 解的是\textbf{反馈 Nash}均衡——策略是状态的函数 $\mathbf u_i=-\mathbf K_i(t)\mathbf x$。还有另一种解概念叫\textbf{开环 Nash}：策略是\emph{时间}的函数 $\mathbf u_i(t)$，双方在 $t=0$ 各自承诺一条完整计划、之后不再看状态。这俩在博弈里\textbf{给出不同的均衡、满足不同的方程}——这是单人最优控制里没有的现象。回忆单人确定性最优控制：开环最优解（PMP 得到的 $u^*(t)$）和反馈最优解（HJB 得到的 $u^*(x,t)$）沿最优轨迹\emph{完全重合}——确定性、无对手时，“先想好一条路”和“边走边看”结果一样，所以平时不区分。但一旦进入博弈，区别就实质化了：均衡依赖\emph{每个玩家掌握什么信息}。开环 Nash 假设双方只在开局看一眼初始状态、之后各按计划走（信息少）；反馈 Nash 假设双方时刻都能观测状态、随时调整（信息多）。信息不同，最优响应不同，于是均衡不同——这正是 \cref{sec:dg-isaacs} 反复强调的“博弈的解依赖信息结构”在 LQ 世界里的具体化身。\textbf{反馈 Nash}：解耦合 Riccati \cref{eq:dg-coupled}，得反馈增益 $\mathbf K_i(t)$，本章与后续 iLQGames 用的就是它；\textbf{开环 Nash}：解一组耦合的两点边值问题（PMP 式协态方程联立），得开环输入 $\mathbf u_i(t)$，方程结构与 \cref{eq:dg-coupled} 不同。

\textbf{本章为什么用反馈 Nash？}\;三个理由：其一，反馈策略有扰动拒绝能力（\cref{sec:dg-isaacs} 已述），机器人实际部署的就是“看状态实时算控制”的反馈律；其二，iLQGames 的输出天然是反馈增益序列，与 \cref{eq:dg-coupled} 对接；其三，反馈 Nash 在多数交互场景里是更现实的建模——交通中没人真会“开局承诺一条路、之后闭眼开”。所以后续一律默认反馈 Nash；遇到“开环”字样要警觉，它是另一套方程、另一族均衡。

\begin{pitfall}[概念误区：把多人 LQ 当成各自独立的 LQR]
以为“每人线性二次，那就每人各解一个标准 LQR、各取各的增益”。\textbf{后果}：得到的增益组合不是纳什均衡——当别人也用最优反馈时，你的“最优”并非最优响应；仿真里表现为策略不自洽、系统行为与理论预测对不上。\textbf{根因}：玩家 $i$ 面对的系统矩阵 $\tilde{\mathbf A}_i=\mathbf A-\sum_{j\ne i}\mathbf B_j\mathbf K_j$ 含着别人的增益，必须把所有最优性条件联立，才得到 \cref{eq:dg-coupled} 的耦合方程；独立 LQR 等于扔掉了那些交叉项。\textbf{正解}：联立求解耦合 Riccati \cref{eq:dg-coupled}（同时反向积分 $N$ 条方程）。检验法：算出各 $\mathbf K_i$ 后，固定别人、单独重解玩家 $i$ 的 LQR，看是否得回同一个 $\mathbf K_i$。
\end{pitfall}

\begin{pitfall}[思维陷阱：默认耦合 Riccati 一定有解]
套用单人 LQR“可镇定就有解”的经验，假设耦合 Riccati 也总能解出有界的 $\mathbf P_i$。\textbf{后果}：反向积分时 $\mathbf P_i$ 在有限时间内发散（finite escape），数值爆掉或得到无意义的巨大增益，却找不到原因。\textbf{根因}：交叉项 $\mathbf P_j\mathbf S_j\mathbf P_i$ 破坏了单人 Riccati 的良好性质；即便每个玩家单独的 LQR 都有解，耦合系统仍可能无界解。存在性需要额外条件，不是默认。\textbf{正解}：求解后检查 $\mathbf P_i$ 是否在整个时域上有界、是否半正定；遇到发散时缩短时域或调整权重。实时博弈求解器用正则化、信赖域等手段处理这一不稳定性。
\end{pitfall}

\begin{practice}
\begin{enumerate}
  \item （实现）用任一语言实现两人 LQ 博弈的耦合 Riccati 反向积分器：给定 $\mathbf A,\mathbf B_1,\mathbf B_2,\mathbf Q_1,\mathbf Q_2,\mathbf R_{11},\mathbf R_{22}$ 与终端 $\mathbf Q_i^f$，从 $t=T$ 反向积分 \cref{eq:dg-coupled} 解出 $\mathbf P_1(t),\mathbf P_2(t)$，再由 \cref{eq:dg-gain} 算反馈增益。在一个 2D 双车 unicycle 线性化模型上验证：用得到的反馈做闭环 rollout，确认两车都无法通过单方面改增益降低自己的代价（即停在纳什均衡）。这是后续完整 iLQGames 的第一块积木。
  \item （手推 + 对照）把 \cref{ex:dg-scalar-lq} 的对称两人标量例子改成\emph{非对称}：$s_1\ne s_2$ 或 $q_1\ne q_2$。写出两条耦合代数方程，讨论解的存在性（什么参数下判别式失效、解不存在？），并与“两人各自的单人 LQR 解”数值对照，量化耦合带来的差异。
  \item （概念，连接 iLQGames）解释为什么“任何光滑微分博弈都能局部近似成 LQ 博弈”，以及 iLQGames 每步为什么要解一次本节的耦合 Riccati。如果某一步的耦合 Riccati 无解，iLQGames 该怎么办？（提示：想想正则化 / 信赖域。）
\end{enumerate}
\end{practice}

% ════════════════════════════════════════════════════════════════
\section{维度诅咒与三条绕行路}
\label{sec:dg-curse}

\paragraph{动机：6 维天花板与真实世界的落差。}
\cref{sec:dg-numerics} 的工程权衡表里那条指数增长的“维数”是 HJI 直接应用的死穴：网格点数随状态维数 $d$ 按 $N^d$ 爆炸，$N=50$、$d=6$ 已是 $50^6\approx1.5\times10^{10}$ 个点的单机极限。可真实问题动辄超过 6 维：两架四旋翼的相对追逃就是 6 维起步，三机就 12 维以上；一个带姿态的固定翼是 6--8 维；多机协调更是维数随机器数线性叠加。所以“HJI 给精确安全证书”这件好事，必须配上能绕过维度诅咒的办法，否则只能停在玩具系统上。下面三条路代表三种不同的取舍哲学。

\paragraph{路一：状态解耦与序贯轨迹规划（STP）。}
\textbf{思想}：很多高维系统的状态之间耦合是\emph{稀疏}的——某些子状态只通过少数共享变量与其余部分相连。Chen 等（2018）证明：对一类满足解耦结构的非线性系统，可以把高维可达集计算\emph{分解}成若干低维子系统的计算，子系统之间只通过共享的状态/控制/干扰耦合\cite{chen2018decomposition}。在多机追逃/规划里，这条路具体化为\textbf{序贯轨迹规划（Sequential Trajectory Planning, STP）}（含对抗入侵者下的鲁棒版本\cite{chen2017robust}）：给多架飞行器排一个优先级，按优先级逐个规划，每架只需在自己的约 6 维相对状态空间里解一个 HJI 子问题，把已规划的高优先级飞行器当作要避让的动态障碍。这样 $N$ 架飞行器的 $6N$ 维联合问题，被拆成 $N$ 个各约 6 维的子问题，计算量从指数降到随机器数线性叠加。\textbf{取舍}：在可解耦/可定优先级的结构下，STP 给出的安全保证仍是\emph{严格}的（不是近似）。代价是它\emph{要求系统具备这种解耦结构或可接受优先级排序}——对状态强耦合、无法分优先级的问题（如真正同时博弈、对称竞速），解耦不适用或会引入保守性。

\paragraph{路二：FaSTrack——把规划与安全分离。}
\textbf{思想}：FaSTrack（Herbert 等 2017）换了个角度：不直接对全维系统算 HJI，而是把问题拆成\emph{两层}\cite{herbert2017fastrack}。\textbf{规划层}：用一个低维、快速的简化模型实时规划（比如把无人机当成匀速运动的点）——快但不精确。\textbf{追踪层}：用 HJI \emph{离线预计算}一个\textbf{追踪误差界（Tracking Error Bound, TEB）}——把“简化规划模型”与“全维真实系统”之间的相对动力学建成一个零和追逃博弈（真实系统追规划点、最坏情况干扰当逃跑者），算出真实系统跟踪规划轨迹时误差不会超过的那个有保证的界（一个不变的“误差球/管”）。在线运行时，规划层随便用什么快速规划器都行，只要把规划出的轨迹按 TEB 向外膨胀（给障碍留出误差球的余量），就\emph{保证}真实系统全程安全——HJI 只在离线那一步、且只在低维的相对系统上算一次。\textbf{取舍}：FaSTrack 给\emph{严格安全保证}且\emph{在线实时}，是工程上很受欢迎的折中。代价是：TEB 需要离线预计算（限制了能处理的相对系统维数），且误差界往往\emph{偏保守}（为覆盖最坏情况，膨胀量可能过大，挤掉可行空间）。

\begin{note}
\textbf{多视角理解（FaSTrack 的 TEB $\leftrightarrow$ \cref{sec:dg-reach} 的 BRT / Tube MPC 的 RPI）.}\;FaSTrack 的追踪误差界，本质又是一根“管”——和 BRT、Tube MPC 的 RPI 集是同一族思想。相似之处：都把最坏情况建成对手、算一个有保证的不变/可达管，把“安全”归结为“待在管内”。不同之处：FaSTrack 算的是“真实系统相对规划模型”的误差管，并把它和一个独立的快速规划器解耦——这就是它能实时的关键。这条把 \cref{sec:dg-reach}、\cref{sec:dg-lq}（局部近似）、FaSTrack 串起来的线索再次印证：本章反复在做同一件事——用对抗博弈算一根保证安全的管/集，区别只在算的是什么、怎么让它可算。
\end{note}

\paragraph{路三：DeepReach——用神经网络近似值函数。}
\textbf{思想}：前两条路都还在“精确求解”的框架内（只是巧妙地降维或分层）。DeepReach（Bansal 与 Tomlin，2021）则彻底换赛道：\emph{放弃网格，用神经网络直接近似 HJI 值函数} $V(x,t)$\cite{bansal2021deepreach}。它的关键技术有两点：其一，用\emph{正弦激活的神经网络}（周期激活，源自隐式神经表示 SIREN 的思路）当函数逼近器——这类网络擅长拟合带高频细节（如棱）的函数；其二，\emph{自监督训练}：不需要标注数据，直接把 HJI 方程的残差（以及边界/终端条件的误差）当作损失函数，让网络去最小化它，于是网络收敛时就近似满足了 HJI 方程。由于神经网络的参数量不随状态维数指数增长，DeepReach 能处理 10 维以上的系统，突破了网格法的 6 维天花板。

\textbf{把“自监督”拆慢一点：损失函数到底是什么。}\;设神经网络 $V_\theta(x,t)$（参数 $\theta$）是对值函数的猜测。我们不知道真值 $V$ 长什么样（没法监督），但我们\emph{知道真值必须满足的方程}——HJI \cref{eq:dg-brs-pde}。于是不去拟合“标签”，而是去逼网络满足那条方程，分两块罚：\textbf{PDE 残差项}——在状态-时间空间里随机采一大批点 $(x_k,t_k)$，每个点上算 HJI 方程左端 $|D_t V_\theta+H(x_k,\nabla V_\theta)|^2$，真解这一项为零，把它当损失去最小化就是在逼网络处处满足方程（其中 $\nabla V_\theta,D_t V_\theta$ 都用自动微分对网络直接求出，这正是正弦激活网络的用武之地）；\textbf{边界/终端项}——在 $t=0$ 处罚 $|V_\theta(x,0)-g(x)|^2$，把终端条件钉住。把两项加权求和当总损失、用梯度下降训练 $\theta$。“自监督”的含义就在这里：监督信号不是外部标签，而是方程本身。这套“拿 PDE 残差当损失”的思路就是物理信息神经网络（PINN）的范式，DeepReach 是它在 HJI 可达性上的落地。直觉上，网格法是在每个格点上\emph{精确}解方程（故受 $N^d$ 之困），DeepReach 是用一个连续函数 $V_\theta$ 去\emph{近似}满足方程（故能上高维，但只是近似）。\textbf{取舍}：代价是\emph{丢掉了全局最优/安全保证}。网格法在收敛时给的是有保证的解，而 DeepReach 是\emph{近似}——网络可能在某些区域偏离真值，导致用它做安全验证时出现违反：在 DeepReach 的形式化验证后续工作\cite{lin2023generating}里，直接用基线 DeepReach 值函数做的某 reach--avoid rollout 仅约 $70\%$（$250$ 条中约 $70\%$）满足安全约束，这正是“近似 $\ne$ 保证”的实证。所以 DeepReach 是“用保证换可扩展”，需要配合后验证/认证手段使用。

\begin{table}[H]\centering\small
\caption{绕维度诅咒三条路：按“精确性 vs 可扩展性”排开}
\label{tab:dg-detour}
\begin{tabular}{p{0.16\textwidth} p{0.26\textwidth} p{0.14\textwidth} p{0.12\textwidth} p{0.18\textwidth}}
\toprule
路线 & 怎么绕 & 安全保证 & 可扩展性 & 前提/代价 \\
\midrule
状态解耦 / STP & 把高维拆成低维子问题，只经共享变量耦合 & \textbf{严格}（可解耦下） & 高（线性叠加） & 要求解耦结构 / 可定优先级 \\
FaSTrack & 低维规划 + 离线 HJI 算追踪误差界 & \textbf{严格}（保证安全） & 中--高（在线实时） & 离线预计算 + 误差界偏保守 \\
DeepReach & 神经网络自监督近似值函数 & \emph{近似}（可能违反） & 高（10D+） & 失全局保证，需后验证 \\
\bottomrule
\end{tabular}
\end{table}

\begin{insight}
\textbf{保证的本质是“可控的误差”.}\;网格法和神经近似的根本差别，不在精度高低，而在\emph{误差是否可控、可知}。网格法的误差有理论上界、随分辨率单调收敛，所以它给的是“保证”；神经近似的误差分布未知、无处不在又无处可指，所以它只给“近似”。安全关键系统要的从来不是“平均很准”，而是“最坏不差过某条线”——这条线能不能画出来，就是“保证”与“近似”的分水岭。一条清晰的取舍：\textbf{越想保留严格保证，越要依赖系统结构或离线代价（解耦、FaSTrack）；越想纯粹地扩展到高维，越要放弃保证（DeepReach）}。没有免费的午餐——这也正是 \cref{sec:dg-lq} 那条“退而用局部 LQ 近似”之所以另辟蹊径的原因：它不追求全局可达集，只在每个时刻求一个精确的局部反馈 Nash，用实时性换全局性，是第四种取舍。
\end{insight}

\paragraph{前沿工作与开放问题。}
\textbf{在线更新的值函数}：让值函数以系统参数为条件，从而能在线适应变化的控制/干扰边界，而不必重算——朝“安全保证随环境在线更新”迈进。\textbf{值函数即安全滤波器}：把 HJ 值网络当成一个自适应安全滤波器统一部署到多种系统。\textbf{开放问题}：如何在\emph{非零和 general-sum} 博弈里做 HJI（本章可达性都建在零和上）？如何在线更新神经值函数以适应未知/变化的对手？如何在保留可扩展性的同时给神经近似补上可证明的安全认证（post-hoc certification）？最后这个“认证”问题是神经近似路线能否用在安全关键系统上的命门：网格法的安全保证之所以硬，是因为它在每个网格点上显式算出了值、误差可由分辨率控制、可被单调格式的收敛性背书；神经网络在训练样本附近拟合得好，但在样本稀疏处可能任意偏离，而你事先不知道哪里稀疏。目前的认证思路大致两类——\emph{事后采样验证}（在声称安全的区域大量采样、前向 rollout 看是否真安全，但验不了“所有”状态）；\emph{形式化误差界}（用神经网络验证工具给逼近误差找一个可证上界，再把它当成额外“干扰”膨胀进安全裕度，但当前这类工具可扩展性本身有限）。两条路都还不成熟——这就是为什么本章把神经近似定位为“用保证换可扩展”，而非网格法的全面替代。\pz{以上前沿工作（参数条件可达集、统一安全滤波器“One Filter to Deploy Them All”等）较新、临近资料知识边界，具体归属与结论以原文为准，待联网核实。}

\begin{pitfall}[思维陷阱：以为 DeepReach（神经近似）总是优于网格法]
“神经网络能上高维、又快，那以后都用 DeepReach、不用网格法了。”\textbf{后果}：在需要严格安全保证的场合用近似值函数，出现未被察觉的安全违反；把“近似解”当“有保证的解”用，可能酿成事故。\textbf{根因}：网格法在收敛时给有保证的解，DeepReach 是近似，网络在某些区域可能偏离真值——可扩展性是用“丢掉全局保证”换来的。\textbf{正解}：分清需求。低维、要严格保证 $\to$ 网格法 / FaSTrack / STP；高维、可接受近似 + 后验证 $\to$ DeepReach。检验法：用 DeepReach 时，务必对其值函数做后验证（采样 rollout 验安全、或用认证方法），不要直接信任。
\end{pitfall}

\begin{pitfall}[概念误区：以为状态解耦 / STP 对任意高维系统都适用]
“维度高就拆成低维子问题嘛，解耦总能用。”\textbf{后果}：对状态强耦合、无法分优先级的系统（如真正对称的同时博弈）强行解耦，要么根本拆不开，要么引入巨大保守性，得到的“安全集”过小、没法用。\textbf{根因}：解耦/STP 依赖系统的特定结构（子状态间稀疏耦合、可定优先级）；不具备这种结构的系统，分解会失真或失效。\textbf{正解}：先判断系统是否真的可解耦、能否合理排优先级；不行就转向 FaSTrack（若能离线算低维相对系统）或 DeepReach（若可接受近似），或干脆用 \cref{sec:dg-lq} 的局部近似路线。
\end{pitfall}

\begin{practice}
\begin{enumerate}
  \item （系统性对比）给定一个 8 维多机器人安全问题（如 2 架四旋翼相对状态共 8 维），分别论证：用 STP、FaSTrack、DeepReach 各自需要什么前提、能给什么级别的保证、主要代价是什么？你会优先尝试哪条，为什么？
  \item （辨析保证 vs 近似）解释为什么“DeepReach 在 reach--avoid 上的安全满足率不足 100\%”这件事，对安全关键应用是致命的；以及为什么同样的近似误差对“粗略包络估计”这类用途可能无所谓。由此说明“该不该用近似可达性”取决于什么。
  \item （连接全章）本章给出了应对维度诅咒的若干思路：\cref{sec:dg-lq} 的局部 LQ 近似、本节的三条绕行路。把它们统一放到“精确性—可扩展性—实时性”三维权衡里比较，并指出 iLQGames 落在这个权衡空间的什么位置、为什么自动驾驶/竞速这类应用会选它。
\end{enumerate}
\end{practice}

% ════════════════════════════════════════════════════════════════
\section{追逃经典问题：Homicidal Chauffeur 与 Air3D}
\label{sec:dg-chase}

\paragraph{动机：为什么需要标志性问题。}
抽象的 HJI 方程、可达集、数值格式，都需要一个具体问题来检验理解、调试代码、对比方法。微分博弈领域有两个几乎人人都从它们入门的经典问题：Isaacs 1965 提出的 \textbf{Homicidal Chauffeur}\cite{isaacs1965differential}，和 Mitchell--Tomlin 2005 用作标准基准的 \textbf{Air3D}\cite{mitchell2005time}。它们都是追逃博弈、都用相对坐标、都能在低维网格上算出可达集——是把前面所有理论“跑起来”的第一站。

\paragraph{Homicidal Chauffeur：转弯受限的追捕者。}
Isaacs 用一个略带黑色幽默的设定命名了这个问题（“开车的杀手”）：一个\emph{追捕者是辆车}——速度快、但\emph{转弯半径受限}（不能原地掉头）；一个\emph{逃跑者是个行人}——速度慢、但\emph{全向机动}（可瞬间改变方向）。追捕者想把两者相对距离压到某个捕获半径以下（“撞上”），逃跑者想一直躲开。这是一个零和博弈：同一个“最终相对距离/能否捕获”，一方求小、一方求大，正是 \cref{sec:dg-isaacs} 的设定。它的精妙在于“快但笨”对“慢但灵”的对抗：追捕者速度占优，但转弯的笨拙给了灵活的行人可乘之机——行人可以在追捕者来不及转向的时刻侧身闪过。谁占优、从哪些初始相对位形追捕者必胜，正是 HJI 可达性（\cref{sec:dg-reach}）要回答的——捕获集就是那个 BRT。

\paragraph{Isaacs 的奇异面：值函数的“棱”从哪来。}
分析 Homicidal Chauffeur 时，Isaacs 发现最优值函数上存在若干\textbf{奇异面（singular surfaces）}——值函数在那里不光滑或最优策略在那里特殊。本课重点认识三类（\cref{tab:dg-singular}）。

\begin{table}[H]\centering\small
\caption{Isaacs 奇异面三类}
\label{tab:dg-singular}
\begin{tabular}{p{0.18\textwidth} p{0.40\textwidth} p{0.34\textwidth}}
\toprule
奇异面 & 含义 & 与前文的联系 \\
\midrule
\textbf{障碍面（barrier）} & 捕获区与逃脱区的分界面，跨过它结局定性翻转 & 就是 \cref{sec:dg-reach} 可达集的边界 $\partial\mathcal{R}$ \\
\textbf{散射面（dispersal）} & 某一方有\emph{多个等价最优选择}的面（如行人向左闪、向右闪一样好） & 正是 \cref{sec:dg-viscosity} 的“棱”——多条等价最优路径在此打平，值函数不可微 \\
\textbf{通用面（universal）} & 多条最优轨迹\emph{汇入并沿之同行}的面，像最优路径的“高速公路” & 反映最优反馈策略的结构 \\
\bottomrule
\end{tabular}
\end{table}

\begin{insight}
\cref{sec:dg-viscosity} 抽象地说“值函数会因多条等价最优路径而起棱”，Homicidal Chauffeur 的散射面就是这句话最具体的样子——行人站在散射面上时，向左闪和向右闪被追上的结局完全一样，于是值函数在那里有一道折痕、协态（最优方向）发生跳变。Isaacs 的奇异面分类，本质就是在给 HJI 值函数的非光滑结构“分门别类”。这也再次说明 \cref{sec:dg-viscosity} 的黏性解为什么不可或缺：经典可微解根本描述不了带散射面的值函数。
\end{insight}

\paragraph{Air3D：可达性工具的“Hello World”。}
Air3D 是把两架同型飞行器的追逃写成\textbf{三维相对坐标}的标准基准——几乎每个可达性工具的第一个示例都是它。两架飞行器都以固定速度 $v$ 飞、各有有界转弯率。把“逃逸机”放到“追捕机”的机体坐标系里，用相对位置 $(x,y)$ 和相对朝向 $\psi$ 三个量描述。

\begin{derivation}[Air3D 相对动力学（动坐标系求导带牵连项）]
\label{deriv:dg-air3d}
取追捕机（pursuer）为参考，把坐标原点固定在它身上、$x$ 轴指向它的前进方向。两机都以固定速率 $v$ 平移；追捕机转弯率为 $u$、逃逸机（evader）转弯率为 $d$（这就是双方各自的控制）。相对状态 $(x,y)$ 是逃逸机在追捕机体系里的位置，$\psi$ 是逃逸机朝向相对追捕机朝向的夹角。

\textbf{第一步：写出逃逸机相对追捕机的速度（平移部分）。}\;逃逸机以速率 $v$ 沿自身朝向 $\psi$（相对追捕机）运动，在追捕机体系下平移速度分量是 $(v\cos\psi,\,v\sin\psi)$；追捕机自己以速率 $v$ 沿体系 $x$ 轴正向运动，所以“逃逸机相对追捕机”的平移部分要减去追捕机的速度 $(v,0)$，得 $(v\cos\psi-v,\ v\sin\psi)$。

\textbf{第二步：补上参考系旋转的牵连项。}\;追捕机体坐标系本身在以角速度 $u$ 旋转（追捕机在转弯），而我们要的是在这个\emph{转动坐标系}里观察到的相对位置变化率。转动坐标系里的求导公式给出一个牵连项 $-\,\omega\times r$：对平面、角速度 $u$、位置 $(x,y)$，这个叉积的分量是 $(+u\,y,\ -u\,x)$。直觉是：哪怕逃逸机在世界里不动，追捕机一转头，逃逸机在它眼里的坐标也会“转过去”——这个视觉上的移动就是牵连项。

\textbf{第三步：合并。}\;平移部分加牵连部分：
\[
  \dot x=\underbrace{-v+v\cos\psi}_{\text{平移}}+\underbrace{u\,y}_{\text{牵连}},\qquad \dot y=\underbrace{v\sin\psi}_{\text{平移}}-\underbrace{u\,x}_{\text{牵连}}.
\]
\textbf{第四步：相对朝向。}\;$\psi$ 是两机朝向之差，其变化率就是两个转弯率之差：$\dot\psi=d-u$。
\end{derivation}

合起来得 Air3D 相对动力学：
\begin{equation}
  \dot x=-v+v\cos\psi+u\,y,\qquad \dot y=v\sin\psi-u\,x,\qquad \dot\psi=d-u,
  \label{eq:dg-air3d}
\end{equation}
其中 $u$ 是追捕机的转弯率（争取捕获）、$d$ 是逃逸机的转弯率（争取逃脱），二者各在有界区间取值。这个推导里值得记住的是第二步的牵连项 $(uy,-ux)$——凡是把动力学写到一个\emph{随机器人转动}的坐标系里（机体系、相对系），都会冒出这种 $\omega\times r$ 形式的项，漏掉它是相对坐标建模最常见的错误。

\begin{note}
\textbf{本质洞察（相对坐标把“绝对运动”消成了“牵连项”）.}\;从 6 维绝对坐标降到 3 维相对坐标，省掉的恰好是“两机一起平移、一起旋转”这些对结果无影响的自由度。代价是动力学里多出了牵连项 $(uy,-ux)$——绝对坐标里追捕机的转弯本是一条独立的状态方程 $\dot\psi_A=u$，降维后它“藏”进了相对位置的演化里。这是降维的普遍规律：维度省下来不是免费的，被消去的自由度会以“耦合项/牵连项”的形式重新出现在保留的方程里。看懂这一点，就不会奇怪为什么相对动力学比绝对动力学“看起来更复杂”——复杂度没消失，只是换了个地方待着。这与 \cref{ch:lie} 处理相对位姿（把绝对位姿差写成李群上的相对元）的思想同源。
\end{note}

\paragraph{为什么一定要用相对坐标。}
两架飞行器的绝对状态是 $(x_A,y_A,\psi_A,x_B,y_B,\psi_B)$——六维。但追逃的结果只取决于二者的\emph{相对}位形，与它们在世界里的绝对位置无关。把问题改写到相对坐标 $(x,y,\psi)$，维度从 6 维降到 3 维——而 \cref{sec:dg-numerics} 告诉我们，3 维网格能轻松算、6 维已逼近极限。

\begin{note}
\textbf{多视角理解（相对坐标 $\leftrightarrow$ 降维 / 物理对称性）.}\;“用相对坐标”不只是记号方便，它是在利用问题的\emph{平移与旋转对称性}把维度砍掉一半。相似之处：这和物理里“用相对坐标处理两体问题、把质心运动分离出去”是同一种智慧——无关的对称自由度（绝对位置、整体朝向）被消掉。不同之处：两体问题的相对坐标是为了解析可积，这里是为了让 HJI 网格算得动（直接关系到能不能绕过 \cref{sec:dg-numerics,sec:dg-curse} 的维度诅咒）。一句话：处理追逃博弈的标准第一步，永远是先换到相对坐标降维。
\end{note}

\begin{example}[Air3D 的哈密顿量与捕获集]
\label{ex:dg-air3d-ham}
碰撞/捕获集是相对距离小于半径 $R$ 的圆盘 $\mathcal{T}=\{x^2+y^2\le R^2\}$，带符号距离 $g(x,y,\psi)=\sqrt{x^2+y^2}-R$。角色：追捕机想进入 $\mathcal{T}$（争取存在一个捕获策略，对 $u$ 取 $\min$），逃逸机想避开（对所有逃逸策略都要被捕获才算危险，对 $d$ 取 $\max$）。记协态 $\lambda=(\lambda_x,\lambda_y,\lambda_\psi)=\nabla V$，哈密顿量
\[
  H=\max_{d}\min_{u}\ \Big[\lambda_x(-v+v\cos\psi+u y)+\lambda_y(v\sin\psi-u x)+\lambda_\psi(d-u)\Big].
\]
把含 $u$ 和含 $d$ 的项分别收拢（二者在动力学里线性出现、可分离，故 Isaacs 条件成立，\cref{sec:dg-isaacs}）：
\[
  H=\underbrace{\big[-v\lambda_x+v\lambda_x\cos\psi+v\lambda_y\sin\psi\big]}_{\text{不含控制}}+\min_{u}\big[u(\lambda_x y-\lambda_y x-\lambda_\psi)\big]+\max_{d}\big[d\,\lambda_\psi\big].
\]
由于 $u\in[-\bar u,\bar u]$、$d\in[-\bar d,\bar d]$，两个极值都在边界取到：
\[
  \min_u[\,\cdot\,]=-\bar u\,\lvert\lambda_x y-\lambda_y x-\lambda_\psi\rvert,\qquad \max_d[\,\cdot\,]=\bar d\,\lvert\lambda_\psi\rvert.
\]
最优策略是 bang-bang 的（取决于协态的符号）——这与 \cref{ex:dg-1d} 一维追逃 $u^*=\operatorname{sign}(\lambda)$ 的结构如出一辙，也是零和追逃博弈的普遍特征：在控制约束的盒子里，最优控制总打在角上。把这个 $H$ 连同 BRT 冻结项 \cref{eq:dg-brt-mitchell} 喂进 \cref{sec:dg-numerics} 的数值格式，反向积分得到的零下水平集，就是那张标志性的三维捕获集——逃逸机一旦初始位形落在里面，无论怎么机动都难逃被捕。
\end{example}

\paragraph{类型博弈与程度博弈在数值上的统一。}
\cref{sec:dg-why} 提过 Isaacs 区分“类型博弈”（追上没有，布尔）与“程度博弈”（最终距离多少，连续）。Air3D 表面问的是类型博弈（能不能捕获），但水平集方法把它\emph{变成了程度博弈再求解}：值函数 $V(x,\psi,t)$ 是个连续量（到捕获的某种“代价裕度”），最后取它的零下水平集 $\{V\le0\}$ 才还原出布尔的捕获判断。这正是 \cref{sec:dg-reach} 开头“把布尔问题编码成连续值函数”那句话在追逃语境里的完整兑现——\textbf{水平集方法的威力，就在于它用一个连续的程度博弈，迂回解出了原本难算的类型博弈}，这是整个 HJI 可达性领域得以工程化的根基。

\begin{pitfall}[思维陷阱：用绝对坐标建模追逃，导致维度翻倍]
直接用两个智能体各自的绝对状态 $(x_A,y_A,\psi_A,x_B,y_B,\psi_B)$ 去算可达集。\textbf{后果}：状态维数翻倍（如 3+3=6 维），网格点数随之指数暴涨，本可在 3 维轻松算的问题被推到 6 维的计算极限，甚至算不动。\textbf{根因}：追逃结果只依赖相对位形，绝对位置/整体朝向是与结果无关的对称自由度；用绝对坐标等于白白多算了这些冗余维度。\textbf{正解}：先换到相对坐标（利用平移/旋转对称性），把维度降到最低再算。检验法：问“我的状态里有没有‘整体平移/旋转一下结果不变’的自由度？”——有，就该消掉它。
\end{pitfall}

\begin{pitfall}[概念误区：把散射面上的不可微当成数值 bug]
算 Homicidal Chauffeur / Air3D 时看到值函数有折痕、协态跳变，以为是格式出错或网格太粗。\textbf{后果}：误把物理上真实的散射面当噪声去“修”，或反复加密网格试图让它变光滑（徒劳）。\textbf{根因}：散射面是值函数\emph{真实的、本质的}非光滑结构（多条等价最优路径打平，\cref{sec:dg-viscosity}），不是数值假象；黏性解本就允许这种棱。\textbf{正解}：识别出它是奇异面而非 bug，用尊重黏性的格式（\cref{sec:dg-numerics} WENO）正确捕捉它即可。检验法：加密网格后折痕位置稳定、只是更锐利——那就是真实的散射面，不是数值噪声。
\end{pitfall}

\begin{practice}
\begin{enumerate}
  \item （实现）用任一可达性工具实现 Air3D 追逃的后向可达管计算，可视化三维零下水平集（捕获集）。改变逃逸机的速度或转弯率，观察捕获集如何变化（追捕者更强时它怎么扩张？）。
  \item （概念，连接 \cref{sec:dg-viscosity}）解释 Homicidal Chauffeur 的散射面为什么对应值函数的不可微棱：站在散射面上的一方面临什么样的“等价选择”？为什么这会让协态（最优方向）在该面上跳变？这与 eikonal 例子（\cref{sec:dg-viscosity}）里 $1-|x|$ 在 $x=0$ 的尖有何相似？
  \item （建模，相对坐标）对“地面机器人追一个全向移动的目标”这一平面追逃，自己推导相对坐标下的动力学（设追捕者朝向为参考系）。说明你用了哪些对称性把维度从绝对坐标的几维降到了几维。
\end{enumerate}
\end{practice}

% ════════════════════════════════════════════════════════════════
\section*{本章常见误解汇总}
\begin{table}[H]\centering\small
\begin{tabular}{p{0.44\textwidth} p{0.44\textwidth}}
\toprule
\textbf{常见误解} & \textbf{正确理解} \\
\midrule
博弈是“加了对手约束的最优控制” & 博弈求的是\textbf{均衡（不动点）}，对手轨迹是它自己优化的解、依赖你的策略，不能当固定约束（\cref{sec:dg-why}） \\
博弈均衡唯一且必然存在 & 纳什均衡可能多个、可能不存在（纯策略下）、可能不稳定；均衡选择是\textbf{设计决策}而非数学结果（\cref{sec:dg-why}） \\
$\min_u\max_v=\max_v\min_u$ 总成立 & 普适的只有 $\max\min\le\min\max$；等号需 Isaacs 条件（可分离 / 凸-凹）（\cref{prop:dg-maxmin}） \\
零和博弈只有“一个”值，与信息结构无关 & 值依赖信息结构；安全验证须给干扰非预期策略取下值，才得保守可信的安全集（\cref{sec:dg-info}） \\
HJI 值函数处处可微，可直接用 $\nabla V$ & 值函数普遍有不可微的\textbf{棱}，是黏性解而非经典解（\cref{def:dg-viscosity}） \\
“几乎处处满足方程”就是正确解 & a.e. 弱解不唯一；只有黏性解（黏性极限）才是唯一物理解（\cref{thm:dg-evans}） \\
带 $\min[0,\cdot]$ 的式子算的是 BRS & 那个冻结项给的是 \textbf{BRT（管）}；BRS（集）无冻结项。安全验证要 BRT（\cref{eq:dg-brt-mitchell}） \\
哈密顿量 max/min 随便取 & 由角色唯一决定：求存在取 $\min$、要对所有取 $\max$（\cref{sec:dg-reach}） \\
中心差分就能解 HJI & 须用尊重黏性的单调/WENO 格式 + 数值耗散，否则振荡或收敛到假解（\cref{sec:dg-numerics}） \\
多人 LQ = 各自独立解 LQR & 必须联立解\textbf{耦合 Riccati}；交叉项 $\mathbf P_j\mathbf S_j\mathbf P_i$ 正是“博弈”的全部增量（\cref{eq:dg-coupled}） \\
耦合 Riccati 一定有解 & 不保证；交叉项可致 finite escape，即便各自单人 LQR 都有解（\cref{sec:dg-lq}） \\
神经近似（DeepReach）总优于网格 & 它用\textbf{全局保证}换可扩展；近似误差会变成安全违反，须后验证（\cref{sec:dg-curse}） \\
状态解耦对任意系统都适用 & 依赖稀疏耦合/可定优先级的结构，否则失真或过保守（\cref{sec:dg-curse}） \\
追逃用绝对坐标建模 & 用\textbf{相对坐标}利用平移/旋转对称性降维（如 Air3D 6→3 维）（\cref{deriv:dg-air3d}） \\
\bottomrule
\end{tabular}
\end{table}

% ════════════════════════════════════════════════════════════════
\section*{本章小结}
本章回答了开篇那个问题——\textbf{对抗交互下“最优”是什么}。答案是：不再是“最优”，而是“\textbf{均衡}”；零和博弈的均衡由 HJI 方程刻画，它是单人 HJB 把 $\min_u$ 换成 $\min_u\max_v$ 的博弈推广（\cref{sec:dg-isaacs}）。但 HJI 的解一般不可微，须用黏性解才能唯一确定（\cref{sec:dg-viscosity}）；把它用于安全，就是可达性分析——分清 BRS/BRT/reach--avoid、按角色定 max/min（\cref{sec:dg-reach}）；真要算出来得用尊重黏性的 WENO + TVD-RK 网格法（\cref{sec:dg-numerics}）。全章的主线是一条“\textbf{从理论最优到可算近似的步步退让}”：HJI 给全局最优与严格安全证书，却被维度诅咒卡在约 6 维；于是要么用局部 LQ 近似（\cref{sec:dg-lq} 的耦合 Riccati，iLQGames 的内核），要么走绕维度诅咒的三条路（解耦 / FaSTrack / DeepReach，\cref{sec:dg-curse}），各有“精确性—可扩展性—实时性”的取舍；最后用两个经典追逃问题（\cref{sec:dg-chase}）把这一切落地。本章给博弈线打下的两块地基——\textbf{可达性安全证书}与\textbf{耦合 Riccati}——会在后续章节反复取用。

\begin{table}[H]\centering\small
\caption{符号表}
\label{tab:dg-symbols}
\begin{tabularx}{\linewidth}{lLl}
\toprule
符号 & 含义 & 首次出现 \\
\midrule
$x,\ f(x,u,v)$ & 状态、动力学（耦合系统的流场） & \cref{eq:dg-cost} \\
$u,v$ 或 $a,b$ & 两玩家的控制（一般理论用 $u,v$；可达性用 $a$=控制、$b$=干扰） & \cref{eq:dg-cost,eq:dg-brs-ham} \\
$J,\ L,\ g$ & 总代价、运行（瞬时）代价、终端代价 / 目标集带符号距离 & \cref{eq:dg-cost,eq:dg-sdf} \\
$V(x,t)$ & 值函数（博弈的值 / 可达性值函数） & \cref{eq:dg-hji} \\
$D_t V$ & 值函数对时间的偏导 $\partial V/\partial t$（可达性方程沿用文献记号） & \cref{eq:dg-brs-pde} \\
$\lambda=\nabla_x V$ & 协态（值函数对状态的梯度） & \cref{eq:dg-ham} \\
$H(x,\lambda)$ & 哈密顿量 $\min_u\max_v[\lambda\cdot f+L]$（或按角色的 $\max_a\min_b$） & \cref{eq:dg-ham} \\
$V^-,\,V^+$ & 博弈的下值、上值；相等 $\iff$ Isaacs 条件成立 & \cref{eq:dg-lower} \\
$\alpha,\gamma$ & 非预期策略（看到对方当前动作再回应、不预知未来的映射；$\alpha$ 为下值里 $u$ 的策略，$\gamma$ 为可达性里干扰的策略） & \cref{eq:dg-lower,eq:dg-brs-def} \\
$\zeta(s;x,t,a,b)$ & 从 $(x,t)$ 出发、控制/干扰作用下的状态轨迹在时刻 $s$ 的取值 & \cref{eq:dg-brs-def} \\
$\mathcal{T},\ \mathcal{R}(t),\ \mathcal{R}_{\text{tube}}(t)$ & 目标集、后向可达集（BRS）、后向可达管（BRT） & \cref{eq:dg-brs-def,eq:dg-brt-def} \\
$\mathcal{A},\ c(x)$ & 危险集、其带符号距离（reach--avoid 用） & \cref{eq:dg-reachavoid} \\
$\varepsilon\Delta V$ & 消失黏性项（定义黏性解的正则化扩散） & \cref{deriv:dg-vanishing} \\
$\hat H,\ \alpha,\ p^\pm$ & 数值哈密顿量、耗散系数（最大特征速度）、两侧梯度 & \cref{eq:dg-lf} \\
$\mathbf A,\ \mathbf B_i$ & LQ 博弈的系统矩阵、玩家 $i$ 的输入矩阵 & \cref{eq:dg-lq-dyn} \\
$\mathbf Q_i,\ \mathbf R_{ii}$ & 玩家 $i$ 的状态、控制代价权重 & \cref{eq:dg-lq-cost} \\
$\mathbf P_i,\ \mathbf K_i$ & 玩家 $i$ 的 Riccati 解、反馈增益（$\mathbf u_i=-\mathbf K_i\mathbf x$） & \cref{eq:dg-coupled,eq:dg-gain} \\
$\mathbf S_j=\mathbf B_j\mathbf R_{jj}^{-1}\mathbf B_j^\top$ & 耦合 Riccati 中的简记（对称半正定） & \cref{eq:dg-gain} \\
$\psi,\ R$ & Air3D 相对朝向、捕获半径 & \cref{eq:dg-air3d} \\
\bottomrule
\end{tabularx}
\end{table}

\begin{table}[H]\centering\small
\caption{定理速查表}
\label{tab:dg-thm-cheat}
\begin{tabularx}{\linewidth}{lLl}
\toprule
结论 / 公式 & 一句话说明 & 对应 \\
\midrule
先动者吃亏 \cref{eq:dg-maxmin} & $\max_v\min_u H\le\min_u\max_v H$（普适） & \cref{prop:dg-maxmin} \\
Isaacs 条件 \cref{eq:dg-isaacs-cond} & $\min_u\max_v=\max_v\min_u$ $\iff$ 博弈有值（可分离 / 凸-凹时成立） & \cref{sec:dg-isaacs-cond} \\
HJI 方程 \cref{eq:dg-hji} & $-\partial_t V=\min_u\max_v[\nabla V\cdot f+L]$，$V(\cdot,T)=g$ & \cref{thm:dg-hji} \\
值函数 = 唯一黏性解 & Evans--Souganidis 1984；唯一性由比较原理 & \cref{thm:dg-evans} \\
BRS 方程 \cref{eq:dg-brs-pde} & $D_tV+H=0$，$H=\max_a\min_b\lambda\cdot f$，无冻结项 & \cref{sec:dg-brs} \\
max/min 口诀 & 求存在取 $\min$、要对所有取 $\max$ & \cref{sec:dg-reach} \\
BRT 方程 \cref{eq:dg-brt-vi} & \cref{eq:dg-brs-pde} 加冻结项 $\min[0,H]$ 或 $\min\{D_tV+H,\,g-V\}=0$ & \cref{sec:dg-brt} \\
reach--avoid \cref{eq:dg-reachavoid} & 内层 reach 冻结 + 外层 avoid 否决的双障碍变分不等式 & \cref{thm:dg-reachavoid} \\
数值哈密顿量 \cref{eq:dg-lf} & Lax--Friedrichs：中心差分 $-$ 数值耗散，尊重黏性 & \cref{sec:dg-scheme} \\
耦合 Riccati \cref{eq:dg-coupled} & $-\dot{\mathbf P}_i=\mathbf A^\top \mathbf P_i+\mathbf P_i\mathbf A+\mathbf Q_i-\mathbf P_i\mathbf S_i\mathbf P_i-\sum_{j\ne i}(\mathbf P_j\mathbf S_j\mathbf P_i+\mathbf P_i\mathbf S_j\mathbf P_j)$ & \cref{thm:dg-coupled} \\
反馈增益 \cref{eq:dg-gain} & $\mathbf K_i=\mathbf R_{ii}^{-1}\mathbf B_i^\top \mathbf P_i$ & \cref{sec:dg-lq} \\
绕维度诅咒三路 & 解耦/STP（严格·要结构）· FaSTrack（严格·离线+保守）· DeepReach（近似·高维） & \cref{sec:dg-curse} \\
\bottomrule
\end{tabularx}
\end{table}

\begin{table}[H]\centering\small
\caption{知识点总表}
\label{tab:dg-knowledge}
\begin{tabular}{clll}
\toprule
\# & 知识点 & 核心要点 & 难度 \\
\midrule
1 & 三次认知跨越 & 最优→均衡、预测→预测即均衡、代价已知→代价需推断 & \(\bigstar\bigstar\) \\
2 & Isaacs 方程与条件 & 零和博弈的 HJI；$\min\max=\max\min$ 才有值 & \(\bigstar\bigstar\bigstar\) \\
3 & 黏性解 & 值函数有棱，用黏性极限唯一确定弱解 & \(\bigstar\bigstar\bigstar\) \\
4 & BRS / BRT / reach--avoid & 安全验证 = 可达性；角色定 max/min；安全要 BRT & \(\bigstar\bigstar\bigstar\) \\
5 & 水平集数值方法 & WENO + TVD-RK + CFL；数值耗散 = 消失黏性 & \(\bigstar\bigstar\bigstar\) \\
6 & LQ 博弈与耦合 Riccati & 唯一可闭式解的博弈；iLQGames 的内核 & \(\bigstar\bigstar\bigstar\) \\
7 & 维度诅咒与三条绕行路 & 解耦 / FaSTrack / DeepReach 的精确性-可扩展性权衡 & \(\bigstar\bigstar\bigstar\) \\
8 & 追逃经典问题 & Homicidal Chauffeur / Air3D；奇异面；相对坐标降维 & \(\bigstar\bigstar\) \\
\bottomrule
\end{tabular}
\end{table}

% ════════════════════════════════════════════════════════════════
\section*{延伸阅读}
按难度与角色标注，括号内为定位。
\begin{itemize}
  \item \textbf{奠基与系统}：Isaacs，\emph{Differential Games}，Wiley 1965\cite{isaacs1965differential}（原典、思想源头，行文古老，选读追逃章节即可）；Ba\c{s}ar \& Olsder，\emph{Dynamic Noncooperative Game Theory}，SIAM Classics 1999\cite{basar1999dynamic}（静/动态博弈框架、LQ 耦合 Riccati 的权威整理）。
  \item \textbf{可达性（工程入门首选）}：Bansal、Chen、Herbert、Tomlin，“Hamilton-Jacobi Reachability: A Brief Overview and Recent Advances”，CDC 2017\cite{bansal2017hamilton}（现代综述，BRS/BRT/角色讲得最清楚，最佳入门）；Mitchell、Bayen、Tomlin，IEEE TAC 2005\cite{mitchell2005time}（可达性奠基，含数值方法）；Margellos \& Lygeros，IEEE TAC 2011\cite{margellos2011hamilton}（reach--avoid 表述）；Fisac 等，HSCC 2015\cite{fisac2015reach}（时变 reach--avoid）。
  \item \textbf{黏性解理论（严格基础）}：Crandall \& Lions，Trans. AMS 1983\cite{crandall1983viscosity}（黏性解原典）；Evans \& Souganidis，Indiana Univ. Math. J. 1984\cite{evans1984differential}（博弈值 = 唯一黏性解）。
  \item \textbf{绕维度诅咒}：Bansal \& Tomlin，“DeepReach”，ICRA 2021\cite{bansal2021deepreach}（神经近似突破维度）；Herbert 等，“FaSTrack”，CDC 2017\cite{herbert2017fastrack}（规划与安全分离）；Chen 等，IEEE TAC 2018\cite{chen2018decomposition}（状态解耦）；序贯轨迹规划（STP）的鲁棒版本\cite{chen2017robust}。
  \item \textbf{开源工具}：hj\_reachability（JAX，现代、最干净）、helperOC（MATLAB，经典、教学标准，Air3D 的 Hello World）、optimized\_dp（Python/HeteroCL，GPU 加速网格法）。
\end{itemize}

\section*{本章与后续章节的关系}
\begin{table}[H]\centering\small
\begin{tabularx}{\linewidth}{llL}
\toprule
关系 & 章节 & 衔接点 \\
\midrule
直接前置 & 实时博弈求解器（iLQGames） & \cref{eq:dg-coupled} 耦合 Riccati = iLQGames 每步子问题；反馈 Nash（\cref{sec:dg-lq}）是其输出形式 \\
间接前置 & 逆博弈 & 可微博弈用隐函数定理对 Nash 解的 KKT 求导——KKT 即本章一阶最优性的推广 \\
工具复用 & 安全证书与多智能体学习 & HJI 安全值函数 $\leftrightarrow$ CBF（\cref{sec:dg-frs-cbf} 的最小限制滤波器即 CBF 思想）；可达性安全集 \\
跨部对照 & \cref{ch:control} 控制部（CBF / MPC） & CBF/CLF 安全滤波、鲁棒/Tube MPC 的最深推导在控制部专章；本章只给规划用途 \\
跨专题对照 & 不确定性规划（Tube MPC） & BRT $\leftrightarrow$ Tube MPC 的 RPI 集（\cref{sec:dg-frs-cbf} 已建桥） \\
跨章对照 & \cref{ch:lie} 数学·李群 & $\SE(3)$ 上的追逃博弈（相对位姿在李群上，\cref{deriv:dg-air3d} 同源） \\
基础复用 & \cref{ch:planning} 规划导论 & frozen robot 即“预测—然后—规划”级联的失效；博弈把预测与规划统一 \\
\bottomrule
\end{tabularx}
\end{table}

\section*{故障排查手册}
\begin{enumerate}
  \item \textbf{症状}：可达管（BRT）算出来全空 / 方向反了。\textbf{原因}：reach 与 avoid 的 max/min 取反；动力学方向写错。\textbf{排查}：按口诀核对每个玩家是 $\exists$（min）还是 $\forall$（max）；确认是否启用 BRT 冻结项；把干扰强度调 0，看是否退化为单人可达集（\cref{sec:dg-reach}）。
  \item \textbf{症状}：BRT 边界有毛刺 / 数值振荡。\textbf{原因}：用了不带耗散的格式（如中心差分）；网格太粗。\textbf{排查}：换 WENO / 单调格式；加密网格看是否收敛变光滑（\cref{sec:dg-numerics}）。
  \item \textbf{症状}：细化网格后数值发散 / 出现 NaN。\textbf{原因}：违反 CFL 条件。\textbf{排查}：让 $\Delta t\lesssim\Delta x/\alpha$；网格加密时同步减小 $\Delta t$（\cref{sec:dg-numerics}）。
  \item \textbf{症状}：耦合 Riccati 反向积分发散。\textbf{原因}：解不存在 / 无界（finite escape）。\textbf{排查}：检查 $\mathbf Q_i,\mathbf R_{ii}$ 是否合理；缩短时域 $T$；加正则化（\cref{sec:dg-lq}）。
  \item \textbf{症状}：多人 LQ 增益不对 / 不是纳什。\textbf{原因}：当成独立 LQR、漏了交叉项。\textbf{排查}：检查 RHS 是否含 $-\sum_{j\ne i}(\mathbf P_j\mathbf S_j\mathbf P_i+\mathbf P_i\mathbf S_j\mathbf P_j)$；固定他人增益、重解单人 LQR 对照（\cref{sec:dg-lq}）。
  \item \textbf{症状}：Air3D 可达集不对称 / 形状怪。\textbf{原因}：相对坐标变换写错；$\psi$ 维未取周期边界。\textbf{排查}：重核相对动力学推导（\cref{deriv:dg-air3d}）；确认周期边界含角度维。
  \item \textbf{症状}：神经近似（DeepReach）值函数导致安全违反。\textbf{原因}：神经近似误差，失全局保证。\textbf{排查}：对值函数做 rollout 后验证；上认证方法，勿直接信任（\cref{sec:dg-curse}）。
\end{enumerate}
